% ============================================================
%  Dublin — a modernised Beamer template for academic use
%
%  Palette taken from www.ucd.ie, matched to sam-deegan.com.
%  Designed for Quarto -> Beamer, but works in plain LaTeX too:
%      \documentclass[aspectratio=169]{beamer}
%      % ============================================================
%  Dublin — a modernised Beamer template for academic use
%
%  Palette taken from www.ucd.ie, matched to sam-deegan.com.
%  Designed for Quarto -> Beamer, but works in plain LaTeX too:
%      \documentclass[aspectratio=169]{beamer}
%      % ============================================================
%  Dublin — a modernised Beamer template for academic use
%
%  Palette taken from www.ucd.ie, matched to sam-deegan.com.
%  Designed for Quarto -> Beamer, but works in plain LaTeX too:
%      \documentclass[aspectratio=169]{beamer}
%      % ============================================================
%  Dublin — a modernised Beamer template for academic use
%
%  Palette taken from www.ucd.ie, matched to sam-deegan.com.
%  Designed for Quarto -> Beamer, but works in plain LaTeX too:
%      \documentclass[aspectratio=169]{beamer}
%      \input{dublin-theme.tex}
%
%  Engine: pdflatex, xelatex or lualatex.
% ============================================================

\makeatletter

% ---------- packages ----------
\usepackage[T1]{fontenc}
\usepackage[sfdefault,scale=1.0]{plex-sans}   % IBM Plex Sans
\usepackage[scale=0.92]{plex-mono}            % IBM Plex Mono for code
\usepackage{microtype}
\usepackage{xcolor}
\usepackage{tikz}
\usepackage{tcolorbox}
\usepackage{appendixnumberbeamer}
\usepackage{booktabs}
\usepackage{colortbl}
\usepackage{array}
\usepackage{amssymb}
\usepackage{etoolbox}
\usepackage{graphicx}
\hyphenpenalty=10000 \exhyphenpenalty=10000
\tcbuselibrary{skins}
\usetikzlibrary{calc, fadings, shadows.blur}

% ---------- palette ----------
\definecolor{dgNavy}   {HTML}{04204C}   % headings, title slide ground
\definecolor{dgBlue}   {HTML}{0056A4}   % structure, links, rules
\definecolor{dgBlueDk} {HTML}{003C77}
\definecolor{dgBlueLt} {HTML}{9FC4E0}   % level-1 jumps
\definecolor{dgGreen}  {HTML}{61B77C}   % the one green: badge, bars, list level 1
\definecolor{dgInk}    {HTML}{212529}   % body text
\definecolor{dgMuted}  {HTML}{6C757D}   % captions, footer, notes
\definecolor{dgRule}   {HTML}{D8E0E6}   % hairlines
\definecolor{dgWash}   {HTML}{F2F6F9}   % block grounds
\definecolor{dgGround} {HTML}{FFFFFF}
% The one alert colour, and it has exactly one job: marking the term that is
% about to be substituted from one equation into another, and marking where
% it lands. Nothing else on a slide is this colour, so when it appears the
% room knows what it means without being told twice.
\definecolor{dgAlert}   {HTML}{B3261E}   % substitution highlight
\definecolor{dgAlertBg} {HTML}{FBE9E7}   % its tint, for the landed block

\setbeamercolor{background canvas}{bg=dgGround}
\setbeamercolor{normal text}{fg=dgInk}
\setbeamercolor{structure}{fg=dgBlue}
\setbeamercolor{alerted text}{fg=dgGreen}

% ---------- chrome we do not want ----------
\setbeamertemplate{navigation symbols}{}
\setbeamertemplate{blocks}[default]      % flat: no shadows, no rounding
\setbeamersize{text margin left=1.1cm, text margin right=1.1cm}
\usefonttheme{professionalfonts}

% Maths in the text font, so equations sit with the prose rather than
% dropping into Computer Modern. Swap for newtxsf or Fira Math if you
% prefer; both need packages this does not.
\usepackage[italic, defaultmathsizes]{mathastext}

% ---------- frame title ----------
% Left aligned, generous space above, a hairline beneath, and an optional
% second line via \framesubtitle. No coloured bar: the rule does the work.
\setbeamercolor{frametitle}{fg=dgNavy, bg=}
% Body text is the class size, 10pt. The frame title is pinned at an
% absolute 12pt so it does not shrink with the body.
\setbeamerfont{frametitle}{series=\bfseries, size*={12}{14}}
\setbeamerfont{framesubtitle}{series=\mdseries, size=\small}
\setbeamercolor{framesubtitle}{fg=dgMuted}

\setbeamertemplate{frametitle}{%
  \vspace{0.75em}%
  \begin{minipage}{\dimexpr\paperwidth-2.2cm\relax}
    \raggedright
    \usebeamercolor[fg]{frametitle}\usebeamerfont{frametitle}\insertframetitle%
    \ifx\insertframesubtitle\@empty%
    \else{\\[0.15em]\usebeamercolor[fg]{framesubtitle}%
          \usebeamerfont{framesubtitle}\insertframesubtitle}%
    \fi
    \\[0.35em]
    {\color{dgRule}\rule{\linewidth}{0.6pt}}
  \end{minipage}%
  \vspace{-0.55em}
}

% ---------- footer ----------
% One hairline, short title on the left, slide number on the right.
% Suppressed on the title slide: pandoc emits that frame without [plain],
% so the flag below is how we detect it. \titlepage raises it, the footline
% sees it and lowers it again.
% Appendix frames are numbered I, II, III, restarting at the divider, so a
% long appendix cannot make the talk look twice its length in the footer.
% \appendix itself is a document-level command and cannot be called from
% inside a frame, which is where the Quarto route needs it, so the reset is
% its own macro and \appendixframe / \appendixcontent both call it.
\newif\ifdg@appendix \dg@appendixfalse
\newcommand{\dgappendixstart}{%
  \global\dg@appendixtrue
  \global\setcounter{framenumber}{0}%
}

\newif\ifdg@istitle \dg@istitlefalse
\AtBeginDocument{%
  \let\dg@origtitlepage\titlepage
  \renewcommand{\titlepage}{\global\dg@istitletrue\dg@origtitlepage}%
}
\setbeamertemplate{footline}{%
  \ifdg@istitle
    \global\dg@istitlefalse
    \addtocounter{framenumber}{-1}%
  \else
  \hbox{%
    \begin{beamercolorbox}[wd=\paperwidth, ht=2.6ex, dp=1.6ex, leftskip=1.1cm,
                           rightskip=1.1cm]{}%
      % \\ after a rule that exactly fills the measure has no line left to
      % end, so LaTeX sets the optional argument as text. \par is safe.
      {\color{dgRule}\rule{\linewidth}{0.4pt}}\par\nobreak\vspace{0.35em}
      \color{dgMuted}\scriptsize
      \begingroup\hypersetup{linkcolor=dgMuted}%
      \ifx\dg@shorttitle\@empty
        \textcolor{dgMuted}{\insertshorttitle}%
      \else
        \textcolor{dgMuted}{\dg@shorttitle}%
      \fi\hfill
      \ifdg@appendix
        \textcolor{dgMuted}{\Roman{framenumber}}%
      \else
        \textcolor{dgMuted}{\insertframenumber}%
      \fi
      \endgroup
    \end{beamercolorbox}%
  }%
  \fi
}

% ---------- title slide ----------
%  Configure in the preamble (or Quarto header-includes):
%     \coverphoto{dublin.jpg}        optional background photograph
%     \coverscrim{0.42}              navy over the photo, 0-1
%     \coverbadge{Job Market Paper}  optional green badge
%     \coverlogo{ucd-crest.png}      optional mark, top right
%     \coverlogolift                 soft dark lift behind the mark
%  With no photo the cover is flat navy.
\def\dg@coverphoto{}
\def\dg@coverscrim{0.42}
\def\dg@coverbadge{}
\def\dg@coverlogo{}
\def\dg@thanksqr{}
% A long title makes a noisy footer, and pandoc has no short-title variable,
% so the theme takes one of its own: short-title in the front matter.
\def\dg@shorttitle{}
\newcommand{\shorttitle}[1]{\def\dg@shorttitle{#1}}
\newif\ifdg@logolift \dg@logoliftfalse
\newif\ifdg@logoglow \dg@logoglowfalse
\newif\ifdg@logovignette \dg@logovignettefalse
\newif\ifdg@logoshadow \dg@logoshadowfalse

\newcommand{\coverphoto}[1]{\def\dg@coverphoto{#1}}
\newcommand{\coverscrim}[1]{\def\dg@coverscrim{#1}}
\newcommand{\coverbadge}[1]{\def\dg@coverbadge{#1}}
\newcommand{\coverlogo}[1]{\def\dg@coverlogo{#1}}
\newcommand{\thanksqr}[1]{\def\dg@thanksqr{#1}}
\newcommand{\coverlogolift}{\dg@logolifttrue}
\newcommand{\coverlogoglow}{\dg@logoglowtrue}
% Preferred: darken the corner the mark sits in, or drop a shadow under it.
% Neither puts an edge on the logo, so neither reads as a bad cut-out.
\newcommand{\coverlogovignette}{\dg@logovignettetrue}
\newcommand{\coverlogoshadow}{\dg@logoshadowtrue}

% Photo plus scrim, shared by the cover and the closing slide. With no photo
% it is flat navy: the photograph is the decoration, and a deck without one
% is better plain than patterned.
% #1 is extra navy laid on top, for slides whose text needs more contrast.
\newcommand{\dg@coverbg}[1]{%
  \fill[dgNavy] (current page.south west) rectangle (current page.north east);
  \begin{scope}
    \clip (current page.south west) rectangle (current page.north east);
    \ifx\dg@coverphoto\@empty
    \else
      \node[anchor=center, inner sep=0pt] at (current page.center)
        {\includegraphics[width=\paperwidth, height=\paperheight]{\dg@coverphoto}};
      \fill[dgNavy, opacity=\dg@coverscrim]
        (current page.south west) rectangle (current page.north east);
      \ifdim#1pt>0pt
        \fill[dgNavy, opacity=#1]
          (current page.south west) rectangle (current page.north east);
      \fi
    \fi
  \end{scope}%
}

% A long title must not push the badge off the top of the cover, so the title
% size steps down once it needs more than two lines. Measured, not guessed:
% the natural single-line width against the width of the column it sits in.
\newlength{\dg@titlewd}
\newlength{\dg@titlecol}
\newcommand{\dg@titlesize}{%
  \setlength{\dg@titlecol}{0.66\paperwidth}%
  \settowidth{\dg@titlewd}{\bfseries\fontsize{28}{31}\selectfont\inserttitle}%
  \ifdim\dg@titlewd>2.6\dg@titlecol
    \fontsize{20}{23}\selectfont
  \else
    \ifdim\dg@titlewd>1.6\dg@titlecol
      \fontsize{24}{27}\selectfont
    \else
      \fontsize{28}{31}\selectfont
    \fi
  \fi
}

\defbeamertemplate*{title page}{deegan}[1][]{%
  \begin{tikzpicture}[remember picture, overlay]
    \dg@coverbg{0}

    \node[anchor=west, text width=0.66\paperwidth, align=left, inner sep=0pt]
      at ($(current page.west)+(1.4cm,0.2cm)$)
      {%
        \ifx\dg@coverbadge\@empty\else\badge{\dg@coverbadge}\\[1.0em]\fi
        \color{white}\bfseries\dg@titlesize\inserttitle\\[0.5em]
        \ifx\insertsubtitle\@empty\else
          \color{white}\large\mdseries\insertsubtitle\\[1.1em]\fi
        \textcolor{white}{\rule{2.6cm}{1pt}}\\[1.1em]
        \color{white}\small\insertauthor\\[0.15em]
        \ifx\insertinstitute\@empty\else
          \color{dgRule}\footnotesize\insertinstitute\\[0.35em]\fi
        \color{dgRule}\footnotesize\insertdate%
      };

    \ifx\dg@coverlogo\@empty\else
      \ifdg@logovignette
        \begin{scope}
          \clip (current page.south west) rectangle (current page.north east);
          \fill[dgNavy, opacity=0.55,
                path fading=circle with fuzzy edge 20 percent]
            ($(current page.north east)+(-1.4cm,-1.2cm)$) circle (3.6cm);
        \end{scope}
      \fi
      \ifdg@logoglow
        \fill[white, opacity=0.16,
              path fading=circle with fuzzy edge 15 percent]
          ($(current page.north east)+(-1.80cm,-1.63cm)$) circle (0.88cm);
      \fi
      \ifdg@logolift
        \fill[black, opacity=0.28, rounded corners=3pt]
          ($(current page.north east)+(-2.55cm,-0.62cm)$)
          rectangle ($(current page.north east)+(-0.95cm,-2.58cm)$);
      \fi
      \ifdg@logoshadow
        \node[anchor=north east, inner sep=0pt,
              blur shadow={shadow blur radius=7pt, shadow opacity=55,
                           shadow xshift=0pt, shadow yshift=-1pt}]
          at ($(current page.north east)+(-1.3cm,-0.9cm)$)
          {\includegraphics[height=1.45cm]{\dg@coverlogo}};
      \else
        \node[anchor=north east, inner sep=0pt]
          at ($(current page.north east)+(-1.3cm,-0.9cm)$)
          {\includegraphics[height=1.45cm]{\dg@coverlogo}};
      \fi
    \fi
  \end{tikzpicture}%
}

% ---------- contents, as shown on section dividers ----------
\setbeamercolor{section in toc}{fg=dgNavy}
\setbeamerfont{section in toc}{series=\mdseries, size=\small}
\setbeamertemplate{section in toc}{%
  \leavevmode\usebeamerfont{section in toc}%
  {\color{dgGreen}\bfseries\inserttocsectionnumber.}~%
  \textcolor{dgNavy}{\inserttocsection}%
  \par\vspace{0.4em}%
}
\setbeamertemplate{section in toc shaded}[default][35]

% ---------- section dividers ----------
% Number and name, on the wash rather than full navy, so it reads as a
% breath rather than a second title slide.
\AtBeginSection[]{%
  \begingroup
  \setbeamercolor{background canvas}{bg=dgWash}
  \begin{frame}[plain,noframenumbering]
    \vfill
    \begin{columns}[T,onlytextwidth]
      \begin{column}{0.50\textwidth}
        \raggedright
        {\color{dgBlue}\bfseries\small\MakeUppercase{Section \insertsectionnumber}\par}
        \vspace{0.35em}
        {\color{dgNavy}\bfseries\LARGE\insertsectionhead\par}
        \vspace{0.7em}
        {\color{dgGreen}\rule{2.2cm}{2pt}}
      \end{column}
      \begin{column}{0.42\textwidth}
        \raggedright
        \tableofcontents[currentsection, hideallsubsections]
      \end{column}
    \end{columns}
    \vfill
  \end{frame}
  \endgroup
}

% ---------- equations ----------
% Display maths is numbered, like tables and figures. The Lua filter turns
% every display equation into an `equation` environment; these two lines set
% the tag in the palette rather than in plain black, and stop beamer resetting
% the counter at each section.
% (this file is already inside \makeatletter, so @ is a letter here)
\renewcommand{\theequation}{\arabic{equation}}
\renewcommand{\tagform@}[1]{%
  \maketag@@@{\bfseries\color{dgGreen}(\ignorespaces #1\unskip\@@italiccorr)}%
}

% ---------- code ----------
% Pandoc puts a highlighted block in `Shaded` and a plain one in `verbatim`.
% Both are set small, on the same wash ground the callout boxes use, so a
% syntax example on a slide reads as another panel rather than as raw output.
% fancyvrb arrives with pandoc's highlighting, not with beamer, so a plain
% LaTeX deck that never shows code does not have \fvset at all.
\@ifundefined{fvset}{}{\fvset{fontsize=\scriptsize}}
\AtBeginDocument{%
  \@ifundefined{Shaded}{}{%
    \renewenvironment{Shaded}
      {\begin{tcolorbox}[dgbase, colback=dgWash,
          borderline west={3pt}{0pt}{dgRule},
          top=0.55em, bottom=0.55em, left=0.9em, right=0.9em]}
      {\end{tcolorbox}}%
  }%
}
\AtBeginEnvironment{verbatim}{\scriptsize}

% Pandoc ships a syntax theme in colours that are not ours. Every token is
% pulled back onto the palette: ink for code, muted for comments, blue for
% strings and keywords. Code on a slide is read, not parsed.
\newcommand{\dg@tokcolour}[2]{%
  \expandafter\ifx\csname #1\endcsname\relax\else
    \expandafter\renewcommand\csname #1\endcsname[1]{\textcolor{#2}{##1}}%
  \fi
}
\AtBeginDocument{%
  \@for\dg@tok:={DataTypeTok,DecValTok,BaseNTok,FloatTok,ConstantTok,%
    CharTok,SpecialCharTok,OtherTok,FunctionTok,VariableTok,OperatorTok,%
    BuiltInTok,ExtensionTok,AttributeTok,RegionMarkerTok,NormalTok,%
    PreprocessorTok,InformationTok,WarningTok,ErrorTok,AlertTok}\do{%
    \expandafter\dg@tokcolour\expandafter{\dg@tok}{dgInk}%
  }%
  \@for\dg@tok:={CommentTok,DocumentationTok,AnnotationTok,CommentVarTok}\do{%
    \expandafter\dg@tokcolour\expandafter{\dg@tok}{dgMuted}%
  }%
  \@for\dg@tok:={StringTok,VerbatimStringTok,SpecialStringTok,KeywordTok,%
    ControlFlowTok,ImportTok}\do{%
    \expandafter\dg@tokcolour\expandafter{\dg@tok}{dgBlue}%
  }%
}

% ---------- lists ----------
% One colour scale for both kinds of list, so depth reads the same whether
% a list is numbered or not: green, then blue, then muted. All three use the
% palette green itself -- there is no second, darker green.
% Beamer right-aligns the label in its box, so a narrow square sits further
% right than a numeral of the same level. Each marker is set in a box as wide
% as the numeral it has to line up with, flush left, so a bulleted list and a
% numbered list start on the same vertical.
% Each marker is set in a box exactly as wide as the numeral at that level,
% flush left, so the measurement maintains itself if the fonts change.
\newlength{\dg@lblw}
\newcommand{\dg@marker}[2]{%
  \settowidth{\dg@lblw}{\bfseries #1}%
  \makebox[\dg@lblw][l]{#2}%
}
\setbeamertemplate{itemize item}{%
  \dg@marker{1.}{\color{dgGreen}\rule[0.35ex]{0.5em}{0.5em}}}
\setbeamertemplate{itemize subitem}{%
  \dg@marker{a.}{\color{dgBlue}\rule[0.35ex]{0.4em}{0.4em}}}
\setbeamertemplate{itemize subsubitem}{%
  \dg@marker{i.}{\color{dgMuted}\textendash}}
\setbeamertemplate{enumerate item}{\color{dgGreen}\bfseries\insertenumlabel.}
\setbeamertemplate{enumerate subitem}{\color{dgBlue}\insertsubenumlabel.}
\setbeamertemplate{enumerate subsubitem}{\color{dgMuted}\insertsubsubenumlabel.}
\setbeamercolor{itemize/enumerate body}{fg=dgInk}
\setlength{\leftmargini}{1.4em}

% ---------- speaker notes ----------
% Hidden by default, so a ::: {.notes} block never reaches a student deck.
% The `presenter` profile includes _shared/presenter-notes.tex after this
% file and overrides it with "show notes on second screen=right".
\setbeameroption{hide notes}

% ---------- list spacing ----------
% Beamer sets list items almost on top of one another, which reads as a wall
% of text from the back of a lecture theatre. Top-level items get room to
% breathe; nested items get less than half of it, so the hierarchy stays
% visible without the block coming apart.
%
% This has to be done by redefining \@listi and friends. Setting \itemsep
% from inside the environment does not work: etoolbox's \AtBeginEnvironment
% fires before \list runs, and \list then overwrites the length with the
% value from \@listi. Tested both ways; only this one moves the type.
\def\@listi{\leftmargin\leftmargini
  \topsep 0.40em \parsep \z@ \partopsep \z@ \itemsep 0.75em}
\let\@listI\@listi
\def\@listii{\leftmargin\leftmarginii
  \labelwidth\leftmarginii \advance\labelwidth-\labelsep
  \topsep 0.25em \parsep \z@ \itemsep 0.30em}
\def\@listiii{\leftmargin\leftmarginiii
  \labelwidth\leftmarginiii \advance\labelwidth-\labelsep
  \topsep 0.20em \parsep \z@ \itemsep 0.25em}

% ---------- pandoc's \tightlist ----------
% Pandoc puts \tightlist as the first thing inside every list whose items are
% not separated by blank lines in the markdown, and defines it as
%   \setlength{\itemsep}{0pt}\setlength{\parskip}{0pt}
% That runs INSIDE the list, after \list has already read \@listi, so it
% silently overrode everything set above: measured on the Lecture 1.1 render,
% the gap between one bullet and the next was 13.5pt — exactly the line
% spacing within a single bullet, i.e. no separation at all. Only lists
% written with blank lines between items escaped it, which is why those look
% right and ordinary ones do not.
%
% Neutralising it hands the spacing back to \@listi above, so the 0.75em
% actually reaches the slide. Lists written with blank lines are unchanged:
% pandoc does not emit \tightlist for those.
%
% \def rather than \renewcommand, because this file is also usable outside
% pandoc, where \tightlist does not exist. \AtBeginDocument, because
% pandoc's own \providecommand can come after this file in the preamble.
\AtBeginDocument{\def\tightlist{}}

% ---------- blocks ----------
% Flat, with a left accent bar instead of a title bar.
\setbeamercolor{block title}{fg=dgNavy, bg=dgWash}
\setbeamercolor{block body}{fg=dgInk, bg=dgWash}
\setbeamercolor{block title alerted}{fg=white, bg=dgGreen}
\setbeamercolor{block body alerted}{fg=dgInk, bg=dgWash}
\setbeamercolor{block title example}{fg=white, bg=dgBlue}
\setbeamercolor{block body example}{fg=dgInk, bg=dgWash}
\setbeamerfont{block title}{series=\bfseries, size=\small}

% Continued frames (e.g. a long reference list) keep the plain title.
\setbeamertemplate{frametitle continuation}{}

% ---------- tables ----------
% Panelled summary/results tables, full width, booktabs rules, notes block.
%
%   \begin{dgtable}{l rrrrrr}
%     \dgpanel{7}{Panel A: County Outcome Variables}
%     \thead{Variable} & \thead{Mean} & ... \\ \midrule
%     Personal Income Growth (\%) & 2.145 & ... \\
%   \end{dgtable}
%   \dgnotes{Summary statistics for ...}
%
% Everything is black by default; use \hl{} or \hlg{} to mark the one row
% or cell that carries the argument. Rules span the full text width
% because the table is a tabular* stretched with \extracolsep.
\setlength{\aboverulesep}{0pt}\setlength{\belowrulesep}{0pt}
\setlength{\heavyrulewidth}{0.9pt}\setlength{\lightrulewidth}{0.4pt}
\arrayrulecolor{dgInk}
\renewcommand{\arraystretch}{1.18}

% Every table is \footnotesize, however it was written. \dgtable sets its own
% size, but a markdown table goes through pandoc as a longtable or a plain
% tabular and would otherwise come out at full body size -- far too big for a
% slide, since even a small econ table has more columns than body text does.
% An environment is a group, so each hook undoes itself on \end.
\AtBeginEnvironment{longtable}{\footnotesize}
\AtBeginEnvironment{tabular}{\footnotesize}
\AtBeginEnvironment{tabularx}{\footnotesize}
\AtBeginEnvironment{table}{\footnotesize}

\newcommand{\thead}[1]{\textcolor{dgInk}{\bfseries #1}}
% Colour in a table should mean something. \hl marks the row or cell that
% carries the argument; everything else stays black.
\newcommand{\hl}[1]{\textcolor{dgBlue}{\bfseries #1}}
\newcommand{\hlg}[1]{\textcolor{dgGreen}{\bfseries #1}}
% A panel heading. Emits the heading row only -- put a \midrule above it (or
% below it, if the panel opens with its own column headings).
\newcommand{\dgpanel}[2]{%
  \multicolumn{#1}{@{}l}{\textcolor{dgInk}{\bfseries #2}}\\
}
% \begin{dgtable}[\scriptsize]{l rrr} -- optional first argument sets the
% type size, for a table with more rows than \footnotesize will hold.
\newenvironment{dgtable}[2][\footnotesize]{%
  #1%
  \setlength{\tabcolsep}{0pt}%
  \begin{tabular*}{\linewidth}{@{\extracolsep{\fill}}#2@{}}
  \toprule
}{%
  \bottomrule
  \end{tabular*}%
}
% A colour chip, for the brand-reference slide.
\newcommand{\swatch}[1]{%
  \tikz[baseline=-0.35ex]{\fill[#1] (0,0) rectangle (0.75em,0.75em);}\,%
}

% Fixed-effects rows: a check rather than "Yes". Black, like the rest of the
% table -- colour is reserved for \hl and \hlg.
\newcommand{\yes}{\checkmark}
\newcommand{\no}{}

% \dgnotes lives in dublin-components.tex now, with the other callouts,
% because it is one of them: same wash panel, same left bar. It used to be
% a bold italic "Notes:" label floating in the body, which read as a
% typesetting accident next to a .defn or .takeaway box.

% ---------- captions ----------
% No "Figure 1:" label on a slide; the caption is the sentence under the
% exhibit, set muted and one size down, left aligned to the figure.
\setbeamertemplate{caption}[numbered]
\setbeamertemplate{caption label separator}{: }
\setbeamerfont{caption}{size=\footnotesize}
\setbeamerfont{caption name}{size=\footnotesize, series=\bfseries}
\setbeamercolor{caption}{fg=dgInk}
\setbeamercolor{caption name}{fg=dgInk}
% Pandoc loads the caption package for markdown tables. Left-align those
% captions at the text margin, so a markdown table's caption sits where
% \tabcaption puts one rather than centred over a narrower table.
\AtBeginDocument{%
  \@ifpackageloaded{caption}{%
    \captionsetup{justification=raggedright, singlelinecheck=false}%
  }{}%
}

% Pandoc emits \includegraphics[...,height=\textheight] for a markdown
% image, which fills the frame and pushes the caption off the slide.
% A figure environment is a group, so shrinking \textheight inside it
% bounds the image and restores itself on exit.
\AtBeginEnvironment{figure}{\setlength{\textheight}{0.50\textheight}}

% ---------- links ----------
\hypersetup{colorlinks=true, linkcolor=dgBlue, urlcolor=dgBlue,
            citecolor=dgBlue, filecolor=dgBlue}

% ---------- bibliography ----------
\setbeamertemplate{bibliography item}{}
\setbeamercolor{bibliography entry author}{fg=dgNavy}
\setbeamercolor{bibliography entry title}{fg=dgInk}
\setbeamerfont{bibliography entry title}{size=\footnotesize}
\setbeamerfont{bibliography entry author}{size=\footnotesize}

\makeatother

%
%  Engine: pdflatex, xelatex or lualatex.
% ============================================================

\makeatletter

% ---------- packages ----------
\usepackage[T1]{fontenc}
\usepackage[sfdefault,scale=1.0]{plex-sans}   % IBM Plex Sans
\usepackage[scale=0.92]{plex-mono}            % IBM Plex Mono for code
\usepackage{microtype}
\usepackage{xcolor}
\usepackage{tikz}
\usepackage{tcolorbox}
\usepackage{appendixnumberbeamer}
\usepackage{booktabs}
\usepackage{colortbl}
\usepackage{array}
\usepackage{amssymb}
\usepackage{etoolbox}
\usepackage{graphicx}
\hyphenpenalty=10000 \exhyphenpenalty=10000
\tcbuselibrary{skins}
\usetikzlibrary{calc, fadings, shadows.blur}

% ---------- palette ----------
\definecolor{dgNavy}   {HTML}{04204C}   % headings, title slide ground
\definecolor{dgBlue}   {HTML}{0056A4}   % structure, links, rules
\definecolor{dgBlueDk} {HTML}{003C77}
\definecolor{dgBlueLt} {HTML}{9FC4E0}   % level-1 jumps
\definecolor{dgGreen}  {HTML}{61B77C}   % the one green: badge, bars, list level 1
\definecolor{dgInk}    {HTML}{212529}   % body text
\definecolor{dgMuted}  {HTML}{6C757D}   % captions, footer, notes
\definecolor{dgRule}   {HTML}{D8E0E6}   % hairlines
\definecolor{dgWash}   {HTML}{F2F6F9}   % block grounds
\definecolor{dgGround} {HTML}{FFFFFF}
% The one alert colour, and it has exactly one job: marking the term that is
% about to be substituted from one equation into another, and marking where
% it lands. Nothing else on a slide is this colour, so when it appears the
% room knows what it means without being told twice.
\definecolor{dgAlert}   {HTML}{B3261E}   % substitution highlight
\definecolor{dgAlertBg} {HTML}{FBE9E7}   % its tint, for the landed block

\setbeamercolor{background canvas}{bg=dgGround}
\setbeamercolor{normal text}{fg=dgInk}
\setbeamercolor{structure}{fg=dgBlue}
\setbeamercolor{alerted text}{fg=dgGreen}

% ---------- chrome we do not want ----------
\setbeamertemplate{navigation symbols}{}
\setbeamertemplate{blocks}[default]      % flat: no shadows, no rounding
\setbeamersize{text margin left=1.1cm, text margin right=1.1cm}
\usefonttheme{professionalfonts}

% Maths in the text font, so equations sit with the prose rather than
% dropping into Computer Modern. Swap for newtxsf or Fira Math if you
% prefer; both need packages this does not.
\usepackage[italic, defaultmathsizes]{mathastext}

% ---------- frame title ----------
% Left aligned, generous space above, a hairline beneath, and an optional
% second line via \framesubtitle. No coloured bar: the rule does the work.
\setbeamercolor{frametitle}{fg=dgNavy, bg=}
% Body text is the class size, 10pt. The frame title is pinned at an
% absolute 12pt so it does not shrink with the body.
\setbeamerfont{frametitle}{series=\bfseries, size*={12}{14}}
\setbeamerfont{framesubtitle}{series=\mdseries, size=\small}
\setbeamercolor{framesubtitle}{fg=dgMuted}

\setbeamertemplate{frametitle}{%
  \vspace{0.75em}%
  \begin{minipage}{\dimexpr\paperwidth-2.2cm\relax}
    \raggedright
    \usebeamercolor[fg]{frametitle}\usebeamerfont{frametitle}\insertframetitle%
    \ifx\insertframesubtitle\@empty%
    \else{\\[0.15em]\usebeamercolor[fg]{framesubtitle}%
          \usebeamerfont{framesubtitle}\insertframesubtitle}%
    \fi
    \\[0.35em]
    {\color{dgRule}\rule{\linewidth}{0.6pt}}
  \end{minipage}%
  \vspace{-0.55em}
}

% ---------- footer ----------
% One hairline, short title on the left, slide number on the right.
% Suppressed on the title slide: pandoc emits that frame without [plain],
% so the flag below is how we detect it. \titlepage raises it, the footline
% sees it and lowers it again.
% Appendix frames are numbered I, II, III, restarting at the divider, so a
% long appendix cannot make the talk look twice its length in the footer.
% \appendix itself is a document-level command and cannot be called from
% inside a frame, which is where the Quarto route needs it, so the reset is
% its own macro and \appendixframe / \appendixcontent both call it.
\newif\ifdg@appendix \dg@appendixfalse
\newcommand{\dgappendixstart}{%
  \global\dg@appendixtrue
  \global\setcounter{framenumber}{0}%
}

\newif\ifdg@istitle \dg@istitlefalse
\AtBeginDocument{%
  \let\dg@origtitlepage\titlepage
  \renewcommand{\titlepage}{\global\dg@istitletrue\dg@origtitlepage}%
}
\setbeamertemplate{footline}{%
  \ifdg@istitle
    \global\dg@istitlefalse
    \addtocounter{framenumber}{-1}%
  \else
  \hbox{%
    \begin{beamercolorbox}[wd=\paperwidth, ht=2.6ex, dp=1.6ex, leftskip=1.1cm,
                           rightskip=1.1cm]{}%
      % \\ after a rule that exactly fills the measure has no line left to
      % end, so LaTeX sets the optional argument as text. \par is safe.
      {\color{dgRule}\rule{\linewidth}{0.4pt}}\par\nobreak\vspace{0.35em}
      \color{dgMuted}\scriptsize
      \begingroup\hypersetup{linkcolor=dgMuted}%
      \ifx\dg@shorttitle\@empty
        \textcolor{dgMuted}{\insertshorttitle}%
      \else
        \textcolor{dgMuted}{\dg@shorttitle}%
      \fi\hfill
      \ifdg@appendix
        \textcolor{dgMuted}{\Roman{framenumber}}%
      \else
        \textcolor{dgMuted}{\insertframenumber}%
      \fi
      \endgroup
    \end{beamercolorbox}%
  }%
  \fi
}

% ---------- title slide ----------
%  Configure in the preamble (or Quarto header-includes):
%     \coverphoto{dublin.jpg}        optional background photograph
%     \coverscrim{0.42}              navy over the photo, 0-1
%     \coverbadge{Job Market Paper}  optional green badge
%     \coverlogo{ucd-crest.png}      optional mark, top right
%     \coverlogolift                 soft dark lift behind the mark
%  With no photo the cover is flat navy.
\def\dg@coverphoto{}
\def\dg@coverscrim{0.42}
\def\dg@coverbadge{}
\def\dg@coverlogo{}
\def\dg@thanksqr{}
% A long title makes a noisy footer, and pandoc has no short-title variable,
% so the theme takes one of its own: short-title in the front matter.
\def\dg@shorttitle{}
\newcommand{\shorttitle}[1]{\def\dg@shorttitle{#1}}
\newif\ifdg@logolift \dg@logoliftfalse
\newif\ifdg@logoglow \dg@logoglowfalse
\newif\ifdg@logovignette \dg@logovignettefalse
\newif\ifdg@logoshadow \dg@logoshadowfalse

\newcommand{\coverphoto}[1]{\def\dg@coverphoto{#1}}
\newcommand{\coverscrim}[1]{\def\dg@coverscrim{#1}}
\newcommand{\coverbadge}[1]{\def\dg@coverbadge{#1}}
\newcommand{\coverlogo}[1]{\def\dg@coverlogo{#1}}
\newcommand{\thanksqr}[1]{\def\dg@thanksqr{#1}}
\newcommand{\coverlogolift}{\dg@logolifttrue}
\newcommand{\coverlogoglow}{\dg@logoglowtrue}
% Preferred: darken the corner the mark sits in, or drop a shadow under it.
% Neither puts an edge on the logo, so neither reads as a bad cut-out.
\newcommand{\coverlogovignette}{\dg@logovignettetrue}
\newcommand{\coverlogoshadow}{\dg@logoshadowtrue}

% Photo plus scrim, shared by the cover and the closing slide. With no photo
% it is flat navy: the photograph is the decoration, and a deck without one
% is better plain than patterned.
% #1 is extra navy laid on top, for slides whose text needs more contrast.
\newcommand{\dg@coverbg}[1]{%
  \fill[dgNavy] (current page.south west) rectangle (current page.north east);
  \begin{scope}
    \clip (current page.south west) rectangle (current page.north east);
    \ifx\dg@coverphoto\@empty
    \else
      \node[anchor=center, inner sep=0pt] at (current page.center)
        {\includegraphics[width=\paperwidth, height=\paperheight]{\dg@coverphoto}};
      \fill[dgNavy, opacity=\dg@coverscrim]
        (current page.south west) rectangle (current page.north east);
      \ifdim#1pt>0pt
        \fill[dgNavy, opacity=#1]
          (current page.south west) rectangle (current page.north east);
      \fi
    \fi
  \end{scope}%
}

% A long title must not push the badge off the top of the cover, so the title
% size steps down once it needs more than two lines. Measured, not guessed:
% the natural single-line width against the width of the column it sits in.
\newlength{\dg@titlewd}
\newlength{\dg@titlecol}
\newcommand{\dg@titlesize}{%
  \setlength{\dg@titlecol}{0.66\paperwidth}%
  \settowidth{\dg@titlewd}{\bfseries\fontsize{28}{31}\selectfont\inserttitle}%
  \ifdim\dg@titlewd>2.6\dg@titlecol
    \fontsize{20}{23}\selectfont
  \else
    \ifdim\dg@titlewd>1.6\dg@titlecol
      \fontsize{24}{27}\selectfont
    \else
      \fontsize{28}{31}\selectfont
    \fi
  \fi
}

\defbeamertemplate*{title page}{deegan}[1][]{%
  \begin{tikzpicture}[remember picture, overlay]
    \dg@coverbg{0}

    \node[anchor=west, text width=0.66\paperwidth, align=left, inner sep=0pt]
      at ($(current page.west)+(1.4cm,0.2cm)$)
      {%
        \ifx\dg@coverbadge\@empty\else\badge{\dg@coverbadge}\\[1.0em]\fi
        \color{white}\bfseries\dg@titlesize\inserttitle\\[0.5em]
        \ifx\insertsubtitle\@empty\else
          \color{white}\large\mdseries\insertsubtitle\\[1.1em]\fi
        \textcolor{white}{\rule{2.6cm}{1pt}}\\[1.1em]
        \color{white}\small\insertauthor\\[0.15em]
        \ifx\insertinstitute\@empty\else
          \color{dgRule}\footnotesize\insertinstitute\\[0.35em]\fi
        \color{dgRule}\footnotesize\insertdate%
      };

    \ifx\dg@coverlogo\@empty\else
      \ifdg@logovignette
        \begin{scope}
          \clip (current page.south west) rectangle (current page.north east);
          \fill[dgNavy, opacity=0.55,
                path fading=circle with fuzzy edge 20 percent]
            ($(current page.north east)+(-1.4cm,-1.2cm)$) circle (3.6cm);
        \end{scope}
      \fi
      \ifdg@logoglow
        \fill[white, opacity=0.16,
              path fading=circle with fuzzy edge 15 percent]
          ($(current page.north east)+(-1.80cm,-1.63cm)$) circle (0.88cm);
      \fi
      \ifdg@logolift
        \fill[black, opacity=0.28, rounded corners=3pt]
          ($(current page.north east)+(-2.55cm,-0.62cm)$)
          rectangle ($(current page.north east)+(-0.95cm,-2.58cm)$);
      \fi
      \ifdg@logoshadow
        \node[anchor=north east, inner sep=0pt,
              blur shadow={shadow blur radius=7pt, shadow opacity=55,
                           shadow xshift=0pt, shadow yshift=-1pt}]
          at ($(current page.north east)+(-1.3cm,-0.9cm)$)
          {\includegraphics[height=1.45cm]{\dg@coverlogo}};
      \else
        \node[anchor=north east, inner sep=0pt]
          at ($(current page.north east)+(-1.3cm,-0.9cm)$)
          {\includegraphics[height=1.45cm]{\dg@coverlogo}};
      \fi
    \fi
  \end{tikzpicture}%
}

% ---------- contents, as shown on section dividers ----------
\setbeamercolor{section in toc}{fg=dgNavy}
\setbeamerfont{section in toc}{series=\mdseries, size=\small}
\setbeamertemplate{section in toc}{%
  \leavevmode\usebeamerfont{section in toc}%
  {\color{dgGreen}\bfseries\inserttocsectionnumber.}~%
  \textcolor{dgNavy}{\inserttocsection}%
  \par\vspace{0.4em}%
}
\setbeamertemplate{section in toc shaded}[default][35]

% ---------- section dividers ----------
% Number and name, on the wash rather than full navy, so it reads as a
% breath rather than a second title slide.
\AtBeginSection[]{%
  \begingroup
  \setbeamercolor{background canvas}{bg=dgWash}
  \begin{frame}[plain,noframenumbering]
    \vfill
    \begin{columns}[T,onlytextwidth]
      \begin{column}{0.50\textwidth}
        \raggedright
        {\color{dgBlue}\bfseries\small\MakeUppercase{Section \insertsectionnumber}\par}
        \vspace{0.35em}
        {\color{dgNavy}\bfseries\LARGE\insertsectionhead\par}
        \vspace{0.7em}
        {\color{dgGreen}\rule{2.2cm}{2pt}}
      \end{column}
      \begin{column}{0.42\textwidth}
        \raggedright
        \tableofcontents[currentsection, hideallsubsections]
      \end{column}
    \end{columns}
    \vfill
  \end{frame}
  \endgroup
}

% ---------- equations ----------
% Display maths is numbered, like tables and figures. The Lua filter turns
% every display equation into an `equation` environment; these two lines set
% the tag in the palette rather than in plain black, and stop beamer resetting
% the counter at each section.
% (this file is already inside \makeatletter, so @ is a letter here)
\renewcommand{\theequation}{\arabic{equation}}
\renewcommand{\tagform@}[1]{%
  \maketag@@@{\bfseries\color{dgGreen}(\ignorespaces #1\unskip\@@italiccorr)}%
}

% ---------- code ----------
% Pandoc puts a highlighted block in `Shaded` and a plain one in `verbatim`.
% Both are set small, on the same wash ground the callout boxes use, so a
% syntax example on a slide reads as another panel rather than as raw output.
% fancyvrb arrives with pandoc's highlighting, not with beamer, so a plain
% LaTeX deck that never shows code does not have \fvset at all.
\@ifundefined{fvset}{}{\fvset{fontsize=\scriptsize}}
\AtBeginDocument{%
  \@ifundefined{Shaded}{}{%
    \renewenvironment{Shaded}
      {\begin{tcolorbox}[dgbase, colback=dgWash,
          borderline west={3pt}{0pt}{dgRule},
          top=0.55em, bottom=0.55em, left=0.9em, right=0.9em]}
      {\end{tcolorbox}}%
  }%
}
\AtBeginEnvironment{verbatim}{\scriptsize}

% Pandoc ships a syntax theme in colours that are not ours. Every token is
% pulled back onto the palette: ink for code, muted for comments, blue for
% strings and keywords. Code on a slide is read, not parsed.
\newcommand{\dg@tokcolour}[2]{%
  \expandafter\ifx\csname #1\endcsname\relax\else
    \expandafter\renewcommand\csname #1\endcsname[1]{\textcolor{#2}{##1}}%
  \fi
}
\AtBeginDocument{%
  \@for\dg@tok:={DataTypeTok,DecValTok,BaseNTok,FloatTok,ConstantTok,%
    CharTok,SpecialCharTok,OtherTok,FunctionTok,VariableTok,OperatorTok,%
    BuiltInTok,ExtensionTok,AttributeTok,RegionMarkerTok,NormalTok,%
    PreprocessorTok,InformationTok,WarningTok,ErrorTok,AlertTok}\do{%
    \expandafter\dg@tokcolour\expandafter{\dg@tok}{dgInk}%
  }%
  \@for\dg@tok:={CommentTok,DocumentationTok,AnnotationTok,CommentVarTok}\do{%
    \expandafter\dg@tokcolour\expandafter{\dg@tok}{dgMuted}%
  }%
  \@for\dg@tok:={StringTok,VerbatimStringTok,SpecialStringTok,KeywordTok,%
    ControlFlowTok,ImportTok}\do{%
    \expandafter\dg@tokcolour\expandafter{\dg@tok}{dgBlue}%
  }%
}

% ---------- lists ----------
% One colour scale for both kinds of list, so depth reads the same whether
% a list is numbered or not: green, then blue, then muted. All three use the
% palette green itself -- there is no second, darker green.
% Beamer right-aligns the label in its box, so a narrow square sits further
% right than a numeral of the same level. Each marker is set in a box as wide
% as the numeral it has to line up with, flush left, so a bulleted list and a
% numbered list start on the same vertical.
% Each marker is set in a box exactly as wide as the numeral at that level,
% flush left, so the measurement maintains itself if the fonts change.
\newlength{\dg@lblw}
\newcommand{\dg@marker}[2]{%
  \settowidth{\dg@lblw}{\bfseries #1}%
  \makebox[\dg@lblw][l]{#2}%
}
\setbeamertemplate{itemize item}{%
  \dg@marker{1.}{\color{dgGreen}\rule[0.35ex]{0.5em}{0.5em}}}
\setbeamertemplate{itemize subitem}{%
  \dg@marker{a.}{\color{dgBlue}\rule[0.35ex]{0.4em}{0.4em}}}
\setbeamertemplate{itemize subsubitem}{%
  \dg@marker{i.}{\color{dgMuted}\textendash}}
\setbeamertemplate{enumerate item}{\color{dgGreen}\bfseries\insertenumlabel.}
\setbeamertemplate{enumerate subitem}{\color{dgBlue}\insertsubenumlabel.}
\setbeamertemplate{enumerate subsubitem}{\color{dgMuted}\insertsubsubenumlabel.}
\setbeamercolor{itemize/enumerate body}{fg=dgInk}
\setlength{\leftmargini}{1.4em}

% ---------- speaker notes ----------
% Hidden by default, so a ::: {.notes} block never reaches a student deck.
% The `presenter` profile includes _shared/presenter-notes.tex after this
% file and overrides it with "show notes on second screen=right".
\setbeameroption{hide notes}

% ---------- list spacing ----------
% Beamer sets list items almost on top of one another, which reads as a wall
% of text from the back of a lecture theatre. Top-level items get room to
% breathe; nested items get less than half of it, so the hierarchy stays
% visible without the block coming apart.
%
% This has to be done by redefining \@listi and friends. Setting \itemsep
% from inside the environment does not work: etoolbox's \AtBeginEnvironment
% fires before \list runs, and \list then overwrites the length with the
% value from \@listi. Tested both ways; only this one moves the type.
\def\@listi{\leftmargin\leftmargini
  \topsep 0.40em \parsep \z@ \partopsep \z@ \itemsep 0.75em}
\let\@listI\@listi
\def\@listii{\leftmargin\leftmarginii
  \labelwidth\leftmarginii \advance\labelwidth-\labelsep
  \topsep 0.25em \parsep \z@ \itemsep 0.30em}
\def\@listiii{\leftmargin\leftmarginiii
  \labelwidth\leftmarginiii \advance\labelwidth-\labelsep
  \topsep 0.20em \parsep \z@ \itemsep 0.25em}

% ---------- pandoc's \tightlist ----------
% Pandoc puts \tightlist as the first thing inside every list whose items are
% not separated by blank lines in the markdown, and defines it as
%   \setlength{\itemsep}{0pt}\setlength{\parskip}{0pt}
% That runs INSIDE the list, after \list has already read \@listi, so it
% silently overrode everything set above: measured on the Lecture 1.1 render,
% the gap between one bullet and the next was 13.5pt — exactly the line
% spacing within a single bullet, i.e. no separation at all. Only lists
% written with blank lines between items escaped it, which is why those look
% right and ordinary ones do not.
%
% Neutralising it hands the spacing back to \@listi above, so the 0.75em
% actually reaches the slide. Lists written with blank lines are unchanged:
% pandoc does not emit \tightlist for those.
%
% \def rather than \renewcommand, because this file is also usable outside
% pandoc, where \tightlist does not exist. \AtBeginDocument, because
% pandoc's own \providecommand can come after this file in the preamble.
\AtBeginDocument{\def\tightlist{}}

% ---------- blocks ----------
% Flat, with a left accent bar instead of a title bar.
\setbeamercolor{block title}{fg=dgNavy, bg=dgWash}
\setbeamercolor{block body}{fg=dgInk, bg=dgWash}
\setbeamercolor{block title alerted}{fg=white, bg=dgGreen}
\setbeamercolor{block body alerted}{fg=dgInk, bg=dgWash}
\setbeamercolor{block title example}{fg=white, bg=dgBlue}
\setbeamercolor{block body example}{fg=dgInk, bg=dgWash}
\setbeamerfont{block title}{series=\bfseries, size=\small}

% Continued frames (e.g. a long reference list) keep the plain title.
\setbeamertemplate{frametitle continuation}{}

% ---------- tables ----------
% Panelled summary/results tables, full width, booktabs rules, notes block.
%
%   \begin{dgtable}{l rrrrrr}
%     \dgpanel{7}{Panel A: County Outcome Variables}
%     \thead{Variable} & \thead{Mean} & ... \\ \midrule
%     Personal Income Growth (\%) & 2.145 & ... \\
%   \end{dgtable}
%   \dgnotes{Summary statistics for ...}
%
% Everything is black by default; use \hl{} or \hlg{} to mark the one row
% or cell that carries the argument. Rules span the full text width
% because the table is a tabular* stretched with \extracolsep.
\setlength{\aboverulesep}{0pt}\setlength{\belowrulesep}{0pt}
\setlength{\heavyrulewidth}{0.9pt}\setlength{\lightrulewidth}{0.4pt}
\arrayrulecolor{dgInk}
\renewcommand{\arraystretch}{1.18}

% Every table is \footnotesize, however it was written. \dgtable sets its own
% size, but a markdown table goes through pandoc as a longtable or a plain
% tabular and would otherwise come out at full body size -- far too big for a
% slide, since even a small econ table has more columns than body text does.
% An environment is a group, so each hook undoes itself on \end.
\AtBeginEnvironment{longtable}{\footnotesize}
\AtBeginEnvironment{tabular}{\footnotesize}
\AtBeginEnvironment{tabularx}{\footnotesize}
\AtBeginEnvironment{table}{\footnotesize}

\newcommand{\thead}[1]{\textcolor{dgInk}{\bfseries #1}}
% Colour in a table should mean something. \hl marks the row or cell that
% carries the argument; everything else stays black.
\newcommand{\hl}[1]{\textcolor{dgBlue}{\bfseries #1}}
\newcommand{\hlg}[1]{\textcolor{dgGreen}{\bfseries #1}}
% A panel heading. Emits the heading row only -- put a \midrule above it (or
% below it, if the panel opens with its own column headings).
\newcommand{\dgpanel}[2]{%
  \multicolumn{#1}{@{}l}{\textcolor{dgInk}{\bfseries #2}}\\
}
% \begin{dgtable}[\scriptsize]{l rrr} -- optional first argument sets the
% type size, for a table with more rows than \footnotesize will hold.
\newenvironment{dgtable}[2][\footnotesize]{%
  #1%
  \setlength{\tabcolsep}{0pt}%
  \begin{tabular*}{\linewidth}{@{\extracolsep{\fill}}#2@{}}
  \toprule
}{%
  \bottomrule
  \end{tabular*}%
}
% A colour chip, for the brand-reference slide.
\newcommand{\swatch}[1]{%
  \tikz[baseline=-0.35ex]{\fill[#1] (0,0) rectangle (0.75em,0.75em);}\,%
}

% Fixed-effects rows: a check rather than "Yes". Black, like the rest of the
% table -- colour is reserved for \hl and \hlg.
\newcommand{\yes}{\checkmark}
\newcommand{\no}{}

% \dgnotes lives in dublin-components.tex now, with the other callouts,
% because it is one of them: same wash panel, same left bar. It used to be
% a bold italic "Notes:" label floating in the body, which read as a
% typesetting accident next to a .defn or .takeaway box.

% ---------- captions ----------
% No "Figure 1:" label on a slide; the caption is the sentence under the
% exhibit, set muted and one size down, left aligned to the figure.
\setbeamertemplate{caption}[numbered]
\setbeamertemplate{caption label separator}{: }
\setbeamerfont{caption}{size=\footnotesize}
\setbeamerfont{caption name}{size=\footnotesize, series=\bfseries}
\setbeamercolor{caption}{fg=dgInk}
\setbeamercolor{caption name}{fg=dgInk}
% Pandoc loads the caption package for markdown tables. Left-align those
% captions at the text margin, so a markdown table's caption sits where
% \tabcaption puts one rather than centred over a narrower table.
\AtBeginDocument{%
  \@ifpackageloaded{caption}{%
    \captionsetup{justification=raggedright, singlelinecheck=false}%
  }{}%
}

% Pandoc emits \includegraphics[...,height=\textheight] for a markdown
% image, which fills the frame and pushes the caption off the slide.
% A figure environment is a group, so shrinking \textheight inside it
% bounds the image and restores itself on exit.
\AtBeginEnvironment{figure}{\setlength{\textheight}{0.50\textheight}}

% ---------- links ----------
\hypersetup{colorlinks=true, linkcolor=dgBlue, urlcolor=dgBlue,
            citecolor=dgBlue, filecolor=dgBlue}

% ---------- bibliography ----------
\setbeamertemplate{bibliography item}{}
\setbeamercolor{bibliography entry author}{fg=dgNavy}
\setbeamercolor{bibliography entry title}{fg=dgInk}
\setbeamerfont{bibliography entry title}{size=\footnotesize}
\setbeamerfont{bibliography entry author}{size=\footnotesize}

\makeatother

%
%  Engine: pdflatex, xelatex or lualatex.
% ============================================================

\makeatletter

% ---------- packages ----------
\usepackage[T1]{fontenc}
\usepackage[sfdefault,scale=1.0]{plex-sans}   % IBM Plex Sans
\usepackage[scale=0.92]{plex-mono}            % IBM Plex Mono for code
\usepackage{microtype}
\usepackage{xcolor}
\usepackage{tikz}
\usepackage{tcolorbox}
\usepackage{appendixnumberbeamer}
\usepackage{booktabs}
\usepackage{colortbl}
\usepackage{array}
\usepackage{amssymb}
\usepackage{etoolbox}
\usepackage{graphicx}
\hyphenpenalty=10000 \exhyphenpenalty=10000
\tcbuselibrary{skins}
\usetikzlibrary{calc, fadings, shadows.blur}

% ---------- palette ----------
\definecolor{dgNavy}   {HTML}{04204C}   % headings, title slide ground
\definecolor{dgBlue}   {HTML}{0056A4}   % structure, links, rules
\definecolor{dgBlueDk} {HTML}{003C77}
\definecolor{dgBlueLt} {HTML}{9FC4E0}   % level-1 jumps
\definecolor{dgGreen}  {HTML}{61B77C}   % the one green: badge, bars, list level 1
\definecolor{dgInk}    {HTML}{212529}   % body text
\definecolor{dgMuted}  {HTML}{6C757D}   % captions, footer, notes
\definecolor{dgRule}   {HTML}{D8E0E6}   % hairlines
\definecolor{dgWash}   {HTML}{F2F6F9}   % block grounds
\definecolor{dgGround} {HTML}{FFFFFF}
% The one alert colour, and it has exactly one job: marking the term that is
% about to be substituted from one equation into another, and marking where
% it lands. Nothing else on a slide is this colour, so when it appears the
% room knows what it means without being told twice.
\definecolor{dgAlert}   {HTML}{B3261E}   % substitution highlight
\definecolor{dgAlertBg} {HTML}{FBE9E7}   % its tint, for the landed block

\setbeamercolor{background canvas}{bg=dgGround}
\setbeamercolor{normal text}{fg=dgInk}
\setbeamercolor{structure}{fg=dgBlue}
\setbeamercolor{alerted text}{fg=dgGreen}

% ---------- chrome we do not want ----------
\setbeamertemplate{navigation symbols}{}
\setbeamertemplate{blocks}[default]      % flat: no shadows, no rounding
\setbeamersize{text margin left=1.1cm, text margin right=1.1cm}
\usefonttheme{professionalfonts}

% Maths in the text font, so equations sit with the prose rather than
% dropping into Computer Modern. Swap for newtxsf or Fira Math if you
% prefer; both need packages this does not.
\usepackage[italic, defaultmathsizes]{mathastext}

% ---------- frame title ----------
% Left aligned, generous space above, a hairline beneath, and an optional
% second line via \framesubtitle. No coloured bar: the rule does the work.
\setbeamercolor{frametitle}{fg=dgNavy, bg=}
% Body text is the class size, 10pt. The frame title is pinned at an
% absolute 12pt so it does not shrink with the body.
\setbeamerfont{frametitle}{series=\bfseries, size*={12}{14}}
\setbeamerfont{framesubtitle}{series=\mdseries, size=\small}
\setbeamercolor{framesubtitle}{fg=dgMuted}

\setbeamertemplate{frametitle}{%
  \vspace{0.75em}%
  \begin{minipage}{\dimexpr\paperwidth-2.2cm\relax}
    \raggedright
    \usebeamercolor[fg]{frametitle}\usebeamerfont{frametitle}\insertframetitle%
    \ifx\insertframesubtitle\@empty%
    \else{\\[0.15em]\usebeamercolor[fg]{framesubtitle}%
          \usebeamerfont{framesubtitle}\insertframesubtitle}%
    \fi
    \\[0.35em]
    {\color{dgRule}\rule{\linewidth}{0.6pt}}
  \end{minipage}%
  \vspace{-0.55em}
}

% ---------- footer ----------
% One hairline, short title on the left, slide number on the right.
% Suppressed on the title slide: pandoc emits that frame without [plain],
% so the flag below is how we detect it. \titlepage raises it, the footline
% sees it and lowers it again.
% Appendix frames are numbered I, II, III, restarting at the divider, so a
% long appendix cannot make the talk look twice its length in the footer.
% \appendix itself is a document-level command and cannot be called from
% inside a frame, which is where the Quarto route needs it, so the reset is
% its own macro and \appendixframe / \appendixcontent both call it.
\newif\ifdg@appendix \dg@appendixfalse
\newcommand{\dgappendixstart}{%
  \global\dg@appendixtrue
  \global\setcounter{framenumber}{0}%
}

\newif\ifdg@istitle \dg@istitlefalse
\AtBeginDocument{%
  \let\dg@origtitlepage\titlepage
  \renewcommand{\titlepage}{\global\dg@istitletrue\dg@origtitlepage}%
}
\setbeamertemplate{footline}{%
  \ifdg@istitle
    \global\dg@istitlefalse
    \addtocounter{framenumber}{-1}%
  \else
  \hbox{%
    \begin{beamercolorbox}[wd=\paperwidth, ht=2.6ex, dp=1.6ex, leftskip=1.1cm,
                           rightskip=1.1cm]{}%
      % \\ after a rule that exactly fills the measure has no line left to
      % end, so LaTeX sets the optional argument as text. \par is safe.
      {\color{dgRule}\rule{\linewidth}{0.4pt}}\par\nobreak\vspace{0.35em}
      \color{dgMuted}\scriptsize
      \begingroup\hypersetup{linkcolor=dgMuted}%
      \ifx\dg@shorttitle\@empty
        \textcolor{dgMuted}{\insertshorttitle}%
      \else
        \textcolor{dgMuted}{\dg@shorttitle}%
      \fi\hfill
      \ifdg@appendix
        \textcolor{dgMuted}{\Roman{framenumber}}%
      \else
        \textcolor{dgMuted}{\insertframenumber}%
      \fi
      \endgroup
    \end{beamercolorbox}%
  }%
  \fi
}

% ---------- title slide ----------
%  Configure in the preamble (or Quarto header-includes):
%     \coverphoto{dublin.jpg}        optional background photograph
%     \coverscrim{0.42}              navy over the photo, 0-1
%     \coverbadge{Job Market Paper}  optional green badge
%     \coverlogo{ucd-crest.png}      optional mark, top right
%     \coverlogolift                 soft dark lift behind the mark
%  With no photo the cover is flat navy.
\def\dg@coverphoto{}
\def\dg@coverscrim{0.42}
\def\dg@coverbadge{}
\def\dg@coverlogo{}
\def\dg@thanksqr{}
% A long title makes a noisy footer, and pandoc has no short-title variable,
% so the theme takes one of its own: short-title in the front matter.
\def\dg@shorttitle{}
\newcommand{\shorttitle}[1]{\def\dg@shorttitle{#1}}
\newif\ifdg@logolift \dg@logoliftfalse
\newif\ifdg@logoglow \dg@logoglowfalse
\newif\ifdg@logovignette \dg@logovignettefalse
\newif\ifdg@logoshadow \dg@logoshadowfalse

\newcommand{\coverphoto}[1]{\def\dg@coverphoto{#1}}
\newcommand{\coverscrim}[1]{\def\dg@coverscrim{#1}}
\newcommand{\coverbadge}[1]{\def\dg@coverbadge{#1}}
\newcommand{\coverlogo}[1]{\def\dg@coverlogo{#1}}
\newcommand{\thanksqr}[1]{\def\dg@thanksqr{#1}}
\newcommand{\coverlogolift}{\dg@logolifttrue}
\newcommand{\coverlogoglow}{\dg@logoglowtrue}
% Preferred: darken the corner the mark sits in, or drop a shadow under it.
% Neither puts an edge on the logo, so neither reads as a bad cut-out.
\newcommand{\coverlogovignette}{\dg@logovignettetrue}
\newcommand{\coverlogoshadow}{\dg@logoshadowtrue}

% Photo plus scrim, shared by the cover and the closing slide. With no photo
% it is flat navy: the photograph is the decoration, and a deck without one
% is better plain than patterned.
% #1 is extra navy laid on top, for slides whose text needs more contrast.
\newcommand{\dg@coverbg}[1]{%
  \fill[dgNavy] (current page.south west) rectangle (current page.north east);
  \begin{scope}
    \clip (current page.south west) rectangle (current page.north east);
    \ifx\dg@coverphoto\@empty
    \else
      \node[anchor=center, inner sep=0pt] at (current page.center)
        {\includegraphics[width=\paperwidth, height=\paperheight]{\dg@coverphoto}};
      \fill[dgNavy, opacity=\dg@coverscrim]
        (current page.south west) rectangle (current page.north east);
      \ifdim#1pt>0pt
        \fill[dgNavy, opacity=#1]
          (current page.south west) rectangle (current page.north east);
      \fi
    \fi
  \end{scope}%
}

% A long title must not push the badge off the top of the cover, so the title
% size steps down once it needs more than two lines. Measured, not guessed:
% the natural single-line width against the width of the column it sits in.
\newlength{\dg@titlewd}
\newlength{\dg@titlecol}
\newcommand{\dg@titlesize}{%
  \setlength{\dg@titlecol}{0.66\paperwidth}%
  \settowidth{\dg@titlewd}{\bfseries\fontsize{28}{31}\selectfont\inserttitle}%
  \ifdim\dg@titlewd>2.6\dg@titlecol
    \fontsize{20}{23}\selectfont
  \else
    \ifdim\dg@titlewd>1.6\dg@titlecol
      \fontsize{24}{27}\selectfont
    \else
      \fontsize{28}{31}\selectfont
    \fi
  \fi
}

\defbeamertemplate*{title page}{deegan}[1][]{%
  \begin{tikzpicture}[remember picture, overlay]
    \dg@coverbg{0}

    \node[anchor=west, text width=0.66\paperwidth, align=left, inner sep=0pt]
      at ($(current page.west)+(1.4cm,0.2cm)$)
      {%
        \ifx\dg@coverbadge\@empty\else\badge{\dg@coverbadge}\\[1.0em]\fi
        \color{white}\bfseries\dg@titlesize\inserttitle\\[0.5em]
        \ifx\insertsubtitle\@empty\else
          \color{white}\large\mdseries\insertsubtitle\\[1.1em]\fi
        \textcolor{white}{\rule{2.6cm}{1pt}}\\[1.1em]
        \color{white}\small\insertauthor\\[0.15em]
        \ifx\insertinstitute\@empty\else
          \color{dgRule}\footnotesize\insertinstitute\\[0.35em]\fi
        \color{dgRule}\footnotesize\insertdate%
      };

    \ifx\dg@coverlogo\@empty\else
      \ifdg@logovignette
        \begin{scope}
          \clip (current page.south west) rectangle (current page.north east);
          \fill[dgNavy, opacity=0.55,
                path fading=circle with fuzzy edge 20 percent]
            ($(current page.north east)+(-1.4cm,-1.2cm)$) circle (3.6cm);
        \end{scope}
      \fi
      \ifdg@logoglow
        \fill[white, opacity=0.16,
              path fading=circle with fuzzy edge 15 percent]
          ($(current page.north east)+(-1.80cm,-1.63cm)$) circle (0.88cm);
      \fi
      \ifdg@logolift
        \fill[black, opacity=0.28, rounded corners=3pt]
          ($(current page.north east)+(-2.55cm,-0.62cm)$)
          rectangle ($(current page.north east)+(-0.95cm,-2.58cm)$);
      \fi
      \ifdg@logoshadow
        \node[anchor=north east, inner sep=0pt,
              blur shadow={shadow blur radius=7pt, shadow opacity=55,
                           shadow xshift=0pt, shadow yshift=-1pt}]
          at ($(current page.north east)+(-1.3cm,-0.9cm)$)
          {\includegraphics[height=1.45cm]{\dg@coverlogo}};
      \else
        \node[anchor=north east, inner sep=0pt]
          at ($(current page.north east)+(-1.3cm,-0.9cm)$)
          {\includegraphics[height=1.45cm]{\dg@coverlogo}};
      \fi
    \fi
  \end{tikzpicture}%
}

% ---------- contents, as shown on section dividers ----------
\setbeamercolor{section in toc}{fg=dgNavy}
\setbeamerfont{section in toc}{series=\mdseries, size=\small}
\setbeamertemplate{section in toc}{%
  \leavevmode\usebeamerfont{section in toc}%
  {\color{dgGreen}\bfseries\inserttocsectionnumber.}~%
  \textcolor{dgNavy}{\inserttocsection}%
  \par\vspace{0.4em}%
}
\setbeamertemplate{section in toc shaded}[default][35]

% ---------- section dividers ----------
% Number and name, on the wash rather than full navy, so it reads as a
% breath rather than a second title slide.
\AtBeginSection[]{%
  \begingroup
  \setbeamercolor{background canvas}{bg=dgWash}
  \begin{frame}[plain,noframenumbering]
    \vfill
    \begin{columns}[T,onlytextwidth]
      \begin{column}{0.50\textwidth}
        \raggedright
        {\color{dgBlue}\bfseries\small\MakeUppercase{Section \insertsectionnumber}\par}
        \vspace{0.35em}
        {\color{dgNavy}\bfseries\LARGE\insertsectionhead\par}
        \vspace{0.7em}
        {\color{dgGreen}\rule{2.2cm}{2pt}}
      \end{column}
      \begin{column}{0.42\textwidth}
        \raggedright
        \tableofcontents[currentsection, hideallsubsections]
      \end{column}
    \end{columns}
    \vfill
  \end{frame}
  \endgroup
}

% ---------- equations ----------
% Display maths is numbered, like tables and figures. The Lua filter turns
% every display equation into an `equation` environment; these two lines set
% the tag in the palette rather than in plain black, and stop beamer resetting
% the counter at each section.
% (this file is already inside \makeatletter, so @ is a letter here)
\renewcommand{\theequation}{\arabic{equation}}
\renewcommand{\tagform@}[1]{%
  \maketag@@@{\bfseries\color{dgGreen}(\ignorespaces #1\unskip\@@italiccorr)}%
}

% ---------- code ----------
% Pandoc puts a highlighted block in `Shaded` and a plain one in `verbatim`.
% Both are set small, on the same wash ground the callout boxes use, so a
% syntax example on a slide reads as another panel rather than as raw output.
% fancyvrb arrives with pandoc's highlighting, not with beamer, so a plain
% LaTeX deck that never shows code does not have \fvset at all.
\@ifundefined{fvset}{}{\fvset{fontsize=\scriptsize}}
\AtBeginDocument{%
  \@ifundefined{Shaded}{}{%
    \renewenvironment{Shaded}
      {\begin{tcolorbox}[dgbase, colback=dgWash,
          borderline west={3pt}{0pt}{dgRule},
          top=0.55em, bottom=0.55em, left=0.9em, right=0.9em]}
      {\end{tcolorbox}}%
  }%
}
\AtBeginEnvironment{verbatim}{\scriptsize}

% Pandoc ships a syntax theme in colours that are not ours. Every token is
% pulled back onto the palette: ink for code, muted for comments, blue for
% strings and keywords. Code on a slide is read, not parsed.
\newcommand{\dg@tokcolour}[2]{%
  \expandafter\ifx\csname #1\endcsname\relax\else
    \expandafter\renewcommand\csname #1\endcsname[1]{\textcolor{#2}{##1}}%
  \fi
}
\AtBeginDocument{%
  \@for\dg@tok:={DataTypeTok,DecValTok,BaseNTok,FloatTok,ConstantTok,%
    CharTok,SpecialCharTok,OtherTok,FunctionTok,VariableTok,OperatorTok,%
    BuiltInTok,ExtensionTok,AttributeTok,RegionMarkerTok,NormalTok,%
    PreprocessorTok,InformationTok,WarningTok,ErrorTok,AlertTok}\do{%
    \expandafter\dg@tokcolour\expandafter{\dg@tok}{dgInk}%
  }%
  \@for\dg@tok:={CommentTok,DocumentationTok,AnnotationTok,CommentVarTok}\do{%
    \expandafter\dg@tokcolour\expandafter{\dg@tok}{dgMuted}%
  }%
  \@for\dg@tok:={StringTok,VerbatimStringTok,SpecialStringTok,KeywordTok,%
    ControlFlowTok,ImportTok}\do{%
    \expandafter\dg@tokcolour\expandafter{\dg@tok}{dgBlue}%
  }%
}

% ---------- lists ----------
% One colour scale for both kinds of list, so depth reads the same whether
% a list is numbered or not: green, then blue, then muted. All three use the
% palette green itself -- there is no second, darker green.
% Beamer right-aligns the label in its box, so a narrow square sits further
% right than a numeral of the same level. Each marker is set in a box as wide
% as the numeral it has to line up with, flush left, so a bulleted list and a
% numbered list start on the same vertical.
% Each marker is set in a box exactly as wide as the numeral at that level,
% flush left, so the measurement maintains itself if the fonts change.
\newlength{\dg@lblw}
\newcommand{\dg@marker}[2]{%
  \settowidth{\dg@lblw}{\bfseries #1}%
  \makebox[\dg@lblw][l]{#2}%
}
\setbeamertemplate{itemize item}{%
  \dg@marker{1.}{\color{dgGreen}\rule[0.35ex]{0.5em}{0.5em}}}
\setbeamertemplate{itemize subitem}{%
  \dg@marker{a.}{\color{dgBlue}\rule[0.35ex]{0.4em}{0.4em}}}
\setbeamertemplate{itemize subsubitem}{%
  \dg@marker{i.}{\color{dgMuted}\textendash}}
\setbeamertemplate{enumerate item}{\color{dgGreen}\bfseries\insertenumlabel.}
\setbeamertemplate{enumerate subitem}{\color{dgBlue}\insertsubenumlabel.}
\setbeamertemplate{enumerate subsubitem}{\color{dgMuted}\insertsubsubenumlabel.}
\setbeamercolor{itemize/enumerate body}{fg=dgInk}
\setlength{\leftmargini}{1.4em}

% ---------- speaker notes ----------
% Hidden by default, so a ::: {.notes} block never reaches a student deck.
% The `presenter` profile includes _shared/presenter-notes.tex after this
% file and overrides it with "show notes on second screen=right".
\setbeameroption{hide notes}

% ---------- list spacing ----------
% Beamer sets list items almost on top of one another, which reads as a wall
% of text from the back of a lecture theatre. Top-level items get room to
% breathe; nested items get less than half of it, so the hierarchy stays
% visible without the block coming apart.
%
% This has to be done by redefining \@listi and friends. Setting \itemsep
% from inside the environment does not work: etoolbox's \AtBeginEnvironment
% fires before \list runs, and \list then overwrites the length with the
% value from \@listi. Tested both ways; only this one moves the type.
\def\@listi{\leftmargin\leftmargini
  \topsep 0.40em \parsep \z@ \partopsep \z@ \itemsep 0.75em}
\let\@listI\@listi
\def\@listii{\leftmargin\leftmarginii
  \labelwidth\leftmarginii \advance\labelwidth-\labelsep
  \topsep 0.25em \parsep \z@ \itemsep 0.30em}
\def\@listiii{\leftmargin\leftmarginiii
  \labelwidth\leftmarginiii \advance\labelwidth-\labelsep
  \topsep 0.20em \parsep \z@ \itemsep 0.25em}

% ---------- pandoc's \tightlist ----------
% Pandoc puts \tightlist as the first thing inside every list whose items are
% not separated by blank lines in the markdown, and defines it as
%   \setlength{\itemsep}{0pt}\setlength{\parskip}{0pt}
% That runs INSIDE the list, after \list has already read \@listi, so it
% silently overrode everything set above: measured on the Lecture 1.1 render,
% the gap between one bullet and the next was 13.5pt — exactly the line
% spacing within a single bullet, i.e. no separation at all. Only lists
% written with blank lines between items escaped it, which is why those look
% right and ordinary ones do not.
%
% Neutralising it hands the spacing back to \@listi above, so the 0.75em
% actually reaches the slide. Lists written with blank lines are unchanged:
% pandoc does not emit \tightlist for those.
%
% \def rather than \renewcommand, because this file is also usable outside
% pandoc, where \tightlist does not exist. \AtBeginDocument, because
% pandoc's own \providecommand can come after this file in the preamble.
\AtBeginDocument{\def\tightlist{}}

% ---------- blocks ----------
% Flat, with a left accent bar instead of a title bar.
\setbeamercolor{block title}{fg=dgNavy, bg=dgWash}
\setbeamercolor{block body}{fg=dgInk, bg=dgWash}
\setbeamercolor{block title alerted}{fg=white, bg=dgGreen}
\setbeamercolor{block body alerted}{fg=dgInk, bg=dgWash}
\setbeamercolor{block title example}{fg=white, bg=dgBlue}
\setbeamercolor{block body example}{fg=dgInk, bg=dgWash}
\setbeamerfont{block title}{series=\bfseries, size=\small}

% Continued frames (e.g. a long reference list) keep the plain title.
\setbeamertemplate{frametitle continuation}{}

% ---------- tables ----------
% Panelled summary/results tables, full width, booktabs rules, notes block.
%
%   \begin{dgtable}{l rrrrrr}
%     \dgpanel{7}{Panel A: County Outcome Variables}
%     \thead{Variable} & \thead{Mean} & ... \\ \midrule
%     Personal Income Growth (\%) & 2.145 & ... \\
%   \end{dgtable}
%   \dgnotes{Summary statistics for ...}
%
% Everything is black by default; use \hl{} or \hlg{} to mark the one row
% or cell that carries the argument. Rules span the full text width
% because the table is a tabular* stretched with \extracolsep.
\setlength{\aboverulesep}{0pt}\setlength{\belowrulesep}{0pt}
\setlength{\heavyrulewidth}{0.9pt}\setlength{\lightrulewidth}{0.4pt}
\arrayrulecolor{dgInk}
\renewcommand{\arraystretch}{1.18}

% Every table is \footnotesize, however it was written. \dgtable sets its own
% size, but a markdown table goes through pandoc as a longtable or a plain
% tabular and would otherwise come out at full body size -- far too big for a
% slide, since even a small econ table has more columns than body text does.
% An environment is a group, so each hook undoes itself on \end.
\AtBeginEnvironment{longtable}{\footnotesize}
\AtBeginEnvironment{tabular}{\footnotesize}
\AtBeginEnvironment{tabularx}{\footnotesize}
\AtBeginEnvironment{table}{\footnotesize}

\newcommand{\thead}[1]{\textcolor{dgInk}{\bfseries #1}}
% Colour in a table should mean something. \hl marks the row or cell that
% carries the argument; everything else stays black.
\newcommand{\hl}[1]{\textcolor{dgBlue}{\bfseries #1}}
\newcommand{\hlg}[1]{\textcolor{dgGreen}{\bfseries #1}}
% A panel heading. Emits the heading row only -- put a \midrule above it (or
% below it, if the panel opens with its own column headings).
\newcommand{\dgpanel}[2]{%
  \multicolumn{#1}{@{}l}{\textcolor{dgInk}{\bfseries #2}}\\
}
% \begin{dgtable}[\scriptsize]{l rrr} -- optional first argument sets the
% type size, for a table with more rows than \footnotesize will hold.
\newenvironment{dgtable}[2][\footnotesize]{%
  #1%
  \setlength{\tabcolsep}{0pt}%
  \begin{tabular*}{\linewidth}{@{\extracolsep{\fill}}#2@{}}
  \toprule
}{%
  \bottomrule
  \end{tabular*}%
}
% A colour chip, for the brand-reference slide.
\newcommand{\swatch}[1]{%
  \tikz[baseline=-0.35ex]{\fill[#1] (0,0) rectangle (0.75em,0.75em);}\,%
}

% Fixed-effects rows: a check rather than "Yes". Black, like the rest of the
% table -- colour is reserved for \hl and \hlg.
\newcommand{\yes}{\checkmark}
\newcommand{\no}{}

% \dgnotes lives in dublin-components.tex now, with the other callouts,
% because it is one of them: same wash panel, same left bar. It used to be
% a bold italic "Notes:" label floating in the body, which read as a
% typesetting accident next to a .defn or .takeaway box.

% ---------- captions ----------
% No "Figure 1:" label on a slide; the caption is the sentence under the
% exhibit, set muted and one size down, left aligned to the figure.
\setbeamertemplate{caption}[numbered]
\setbeamertemplate{caption label separator}{: }
\setbeamerfont{caption}{size=\footnotesize}
\setbeamerfont{caption name}{size=\footnotesize, series=\bfseries}
\setbeamercolor{caption}{fg=dgInk}
\setbeamercolor{caption name}{fg=dgInk}
% Pandoc loads the caption package for markdown tables. Left-align those
% captions at the text margin, so a markdown table's caption sits where
% \tabcaption puts one rather than centred over a narrower table.
\AtBeginDocument{%
  \@ifpackageloaded{caption}{%
    \captionsetup{justification=raggedright, singlelinecheck=false}%
  }{}%
}

% Pandoc emits \includegraphics[...,height=\textheight] for a markdown
% image, which fills the frame and pushes the caption off the slide.
% A figure environment is a group, so shrinking \textheight inside it
% bounds the image and restores itself on exit.
\AtBeginEnvironment{figure}{\setlength{\textheight}{0.50\textheight}}

% ---------- links ----------
\hypersetup{colorlinks=true, linkcolor=dgBlue, urlcolor=dgBlue,
            citecolor=dgBlue, filecolor=dgBlue}

% ---------- bibliography ----------
\setbeamertemplate{bibliography item}{}
\setbeamercolor{bibliography entry author}{fg=dgNavy}
\setbeamercolor{bibliography entry title}{fg=dgInk}
\setbeamerfont{bibliography entry title}{size=\footnotesize}
\setbeamerfont{bibliography entry author}{size=\footnotesize}

\makeatother

%
%  Engine: pdflatex, xelatex or lualatex.
% ============================================================

\makeatletter

% ---------- packages ----------
\usepackage[T1]{fontenc}
\usepackage[sfdefault,scale=1.0]{plex-sans}   % IBM Plex Sans
\usepackage[scale=0.92]{plex-mono}            % IBM Plex Mono for code
\usepackage{microtype}
\usepackage{xcolor}
\usepackage{tikz}
\usepackage{tcolorbox}
\usepackage{appendixnumberbeamer}
\usepackage{booktabs}
\usepackage{colortbl}
\usepackage{array}
\usepackage{amssymb}
\usepackage{etoolbox}
\usepackage{graphicx}
\hyphenpenalty=10000 \exhyphenpenalty=10000
\tcbuselibrary{skins}
\usetikzlibrary{calc, fadings, shadows.blur}

% ---------- palette ----------
\definecolor{dgNavy}   {HTML}{04204C}   % headings, title slide ground
\definecolor{dgBlue}   {HTML}{0056A4}   % structure, links, rules
\definecolor{dgBlueDk} {HTML}{003C77}
\definecolor{dgBlueLt} {HTML}{9FC4E0}   % level-1 jumps
\definecolor{dgGreen}  {HTML}{61B77C}   % the one green: badge, bars, list level 1
\definecolor{dgInk}    {HTML}{212529}   % body text
\definecolor{dgMuted}  {HTML}{6C757D}   % captions, footer, notes
\definecolor{dgRule}   {HTML}{D8E0E6}   % hairlines
\definecolor{dgWash}   {HTML}{F2F6F9}   % block grounds
\definecolor{dgGround} {HTML}{FFFFFF}
% The one alert colour, and it has exactly one job: marking the term that is
% about to be substituted from one equation into another, and marking where
% it lands. Nothing else on a slide is this colour, so when it appears the
% room knows what it means without being told twice.
\definecolor{dgAlert}   {HTML}{B3261E}   % substitution highlight
\definecolor{dgAlertBg} {HTML}{FBE9E7}   % its tint, for the landed block

\setbeamercolor{background canvas}{bg=dgGround}
\setbeamercolor{normal text}{fg=dgInk}
\setbeamercolor{structure}{fg=dgBlue}
\setbeamercolor{alerted text}{fg=dgGreen}

% ---------- chrome we do not want ----------
\setbeamertemplate{navigation symbols}{}
\setbeamertemplate{blocks}[default]      % flat: no shadows, no rounding
\setbeamersize{text margin left=1.1cm, text margin right=1.1cm}
\usefonttheme{professionalfonts}

% Maths in the text font, so equations sit with the prose rather than
% dropping into Computer Modern. Swap for newtxsf or Fira Math if you
% prefer; both need packages this does not.
\usepackage[italic, defaultmathsizes]{mathastext}

% ---------- frame title ----------
% Left aligned, generous space above, a hairline beneath, and an optional
% second line via \framesubtitle. No coloured bar: the rule does the work.
\setbeamercolor{frametitle}{fg=dgNavy, bg=}
% Body text is the class size, 10pt. The frame title is pinned at an
% absolute 12pt so it does not shrink with the body.
\setbeamerfont{frametitle}{series=\bfseries, size*={12}{14}}
\setbeamerfont{framesubtitle}{series=\mdseries, size=\small}
\setbeamercolor{framesubtitle}{fg=dgMuted}

\setbeamertemplate{frametitle}{%
  \vspace{0.75em}%
  \begin{minipage}{\dimexpr\paperwidth-2.2cm\relax}
    \raggedright
    \usebeamercolor[fg]{frametitle}\usebeamerfont{frametitle}\insertframetitle%
    \ifx\insertframesubtitle\@empty%
    \else{\\[0.15em]\usebeamercolor[fg]{framesubtitle}%
          \usebeamerfont{framesubtitle}\insertframesubtitle}%
    \fi
    \\[0.35em]
    {\color{dgRule}\rule{\linewidth}{0.6pt}}
  \end{minipage}%
  \vspace{-0.55em}
}

% ---------- footer ----------
% One hairline, short title on the left, slide number on the right.
% Suppressed on the title slide: pandoc emits that frame without [plain],
% so the flag below is how we detect it. \titlepage raises it, the footline
% sees it and lowers it again.
% Appendix frames are numbered I, II, III, restarting at the divider, so a
% long appendix cannot make the talk look twice its length in the footer.
% \appendix itself is a document-level command and cannot be called from
% inside a frame, which is where the Quarto route needs it, so the reset is
% its own macro and \appendixframe / \appendixcontent both call it.
\newif\ifdg@appendix \dg@appendixfalse
\newcommand{\dgappendixstart}{%
  \global\dg@appendixtrue
  \global\setcounter{framenumber}{0}%
}

\newif\ifdg@istitle \dg@istitlefalse
\AtBeginDocument{%
  \let\dg@origtitlepage\titlepage
  \renewcommand{\titlepage}{\global\dg@istitletrue\dg@origtitlepage}%
}
\setbeamertemplate{footline}{%
  \ifdg@istitle
    \global\dg@istitlefalse
    \addtocounter{framenumber}{-1}%
  \else
  \hbox{%
    \begin{beamercolorbox}[wd=\paperwidth, ht=2.6ex, dp=1.6ex, leftskip=1.1cm,
                           rightskip=1.1cm]{}%
      % \\ after a rule that exactly fills the measure has no line left to
      % end, so LaTeX sets the optional argument as text. \par is safe.
      {\color{dgRule}\rule{\linewidth}{0.4pt}}\par\nobreak\vspace{0.35em}
      \color{dgMuted}\scriptsize
      \begingroup\hypersetup{linkcolor=dgMuted}%
      \ifx\dg@shorttitle\@empty
        \textcolor{dgMuted}{\insertshorttitle}%
      \else
        \textcolor{dgMuted}{\dg@shorttitle}%
      \fi\hfill
      \ifdg@appendix
        \textcolor{dgMuted}{\Roman{framenumber}}%
      \else
        \textcolor{dgMuted}{\insertframenumber}%
      \fi
      \endgroup
    \end{beamercolorbox}%
  }%
  \fi
}

% ---------- title slide ----------
%  Configure in the preamble (or Quarto header-includes):
%     \coverphoto{dublin.jpg}        optional background photograph
%     \coverscrim{0.42}              navy over the photo, 0-1
%     \coverbadge{Job Market Paper}  optional green badge
%     \coverlogo{ucd-crest.png}      optional mark, top right
%     \coverlogolift                 soft dark lift behind the mark
%  With no photo the cover is flat navy.
\def\dg@coverphoto{}
\def\dg@coverscrim{0.42}
\def\dg@coverbadge{}
\def\dg@coverlogo{}
\def\dg@thanksqr{}
% A long title makes a noisy footer, and pandoc has no short-title variable,
% so the theme takes one of its own: short-title in the front matter.
\def\dg@shorttitle{}
\newcommand{\shorttitle}[1]{\def\dg@shorttitle{#1}}
\newif\ifdg@logolift \dg@logoliftfalse
\newif\ifdg@logoglow \dg@logoglowfalse
\newif\ifdg@logovignette \dg@logovignettefalse
\newif\ifdg@logoshadow \dg@logoshadowfalse

\newcommand{\coverphoto}[1]{\def\dg@coverphoto{#1}}
\newcommand{\coverscrim}[1]{\def\dg@coverscrim{#1}}
\newcommand{\coverbadge}[1]{\def\dg@coverbadge{#1}}
\newcommand{\coverlogo}[1]{\def\dg@coverlogo{#1}}
\newcommand{\thanksqr}[1]{\def\dg@thanksqr{#1}}
\newcommand{\coverlogolift}{\dg@logolifttrue}
\newcommand{\coverlogoglow}{\dg@logoglowtrue}
% Preferred: darken the corner the mark sits in, or drop a shadow under it.
% Neither puts an edge on the logo, so neither reads as a bad cut-out.
\newcommand{\coverlogovignette}{\dg@logovignettetrue}
\newcommand{\coverlogoshadow}{\dg@logoshadowtrue}

% Photo plus scrim, shared by the cover and the closing slide. With no photo
% it is flat navy: the photograph is the decoration, and a deck without one
% is better plain than patterned.
% #1 is extra navy laid on top, for slides whose text needs more contrast.
\newcommand{\dg@coverbg}[1]{%
  \fill[dgNavy] (current page.south west) rectangle (current page.north east);
  \begin{scope}
    \clip (current page.south west) rectangle (current page.north east);
    \ifx\dg@coverphoto\@empty
    \else
      \node[anchor=center, inner sep=0pt] at (current page.center)
        {\includegraphics[width=\paperwidth, height=\paperheight]{\dg@coverphoto}};
      \fill[dgNavy, opacity=\dg@coverscrim]
        (current page.south west) rectangle (current page.north east);
      \ifdim#1pt>0pt
        \fill[dgNavy, opacity=#1]
          (current page.south west) rectangle (current page.north east);
      \fi
    \fi
  \end{scope}%
}

% A long title must not push the badge off the top of the cover, so the title
% size steps down once it needs more than two lines. Measured, not guessed:
% the natural single-line width against the width of the column it sits in.
\newlength{\dg@titlewd}
\newlength{\dg@titlecol}
\newcommand{\dg@titlesize}{%
  \setlength{\dg@titlecol}{0.66\paperwidth}%
  \settowidth{\dg@titlewd}{\bfseries\fontsize{28}{31}\selectfont\inserttitle}%
  \ifdim\dg@titlewd>2.6\dg@titlecol
    \fontsize{20}{23}\selectfont
  \else
    \ifdim\dg@titlewd>1.6\dg@titlecol
      \fontsize{24}{27}\selectfont
    \else
      \fontsize{28}{31}\selectfont
    \fi
  \fi
}

\defbeamertemplate*{title page}{deegan}[1][]{%
  \begin{tikzpicture}[remember picture, overlay]
    \dg@coverbg{0}

    \node[anchor=west, text width=0.66\paperwidth, align=left, inner sep=0pt]
      at ($(current page.west)+(1.4cm,0.2cm)$)
      {%
        \ifx\dg@coverbadge\@empty\else\badge{\dg@coverbadge}\\[1.0em]\fi
        \color{white}\bfseries\dg@titlesize\inserttitle\\[0.5em]
        \ifx\insertsubtitle\@empty\else
          \color{white}\large\mdseries\insertsubtitle\\[1.1em]\fi
        \textcolor{white}{\rule{2.6cm}{1pt}}\\[1.1em]
        \color{white}\small\insertauthor\\[0.15em]
        \ifx\insertinstitute\@empty\else
          \color{dgRule}\footnotesize\insertinstitute\\[0.35em]\fi
        \color{dgRule}\footnotesize\insertdate%
      };

    \ifx\dg@coverlogo\@empty\else
      \ifdg@logovignette
        \begin{scope}
          \clip (current page.south west) rectangle (current page.north east);
          \fill[dgNavy, opacity=0.55,
                path fading=circle with fuzzy edge 20 percent]
            ($(current page.north east)+(-1.4cm,-1.2cm)$) circle (3.6cm);
        \end{scope}
      \fi
      \ifdg@logoglow
        \fill[white, opacity=0.16,
              path fading=circle with fuzzy edge 15 percent]
          ($(current page.north east)+(-1.80cm,-1.63cm)$) circle (0.88cm);
      \fi
      \ifdg@logolift
        \fill[black, opacity=0.28, rounded corners=3pt]
          ($(current page.north east)+(-2.55cm,-0.62cm)$)
          rectangle ($(current page.north east)+(-0.95cm,-2.58cm)$);
      \fi
      \ifdg@logoshadow
        \node[anchor=north east, inner sep=0pt,
              blur shadow={shadow blur radius=7pt, shadow opacity=55,
                           shadow xshift=0pt, shadow yshift=-1pt}]
          at ($(current page.north east)+(-1.3cm,-0.9cm)$)
          {\includegraphics[height=1.45cm]{\dg@coverlogo}};
      \else
        \node[anchor=north east, inner sep=0pt]
          at ($(current page.north east)+(-1.3cm,-0.9cm)$)
          {\includegraphics[height=1.45cm]{\dg@coverlogo}};
      \fi
    \fi
  \end{tikzpicture}%
}

% ---------- contents, as shown on section dividers ----------
\setbeamercolor{section in toc}{fg=dgNavy}
\setbeamerfont{section in toc}{series=\mdseries, size=\small}
\setbeamertemplate{section in toc}{%
  \leavevmode\usebeamerfont{section in toc}%
  {\color{dgGreen}\bfseries\inserttocsectionnumber.}~%
  \textcolor{dgNavy}{\inserttocsection}%
  \par\vspace{0.4em}%
}
\setbeamertemplate{section in toc shaded}[default][35]

% ---------- section dividers ----------
% Number and name, on the wash rather than full navy, so it reads as a
% breath rather than a second title slide.
\AtBeginSection[]{%
  \begingroup
  \setbeamercolor{background canvas}{bg=dgWash}
  \begin{frame}[plain,noframenumbering]
    \vfill
    \begin{columns}[T,onlytextwidth]
      \begin{column}{0.50\textwidth}
        \raggedright
        {\color{dgBlue}\bfseries\small\MakeUppercase{Section \insertsectionnumber}\par}
        \vspace{0.35em}
        {\color{dgNavy}\bfseries\LARGE\insertsectionhead\par}
        \vspace{0.7em}
        {\color{dgGreen}\rule{2.2cm}{2pt}}
      \end{column}
      \begin{column}{0.42\textwidth}
        \raggedright
        \tableofcontents[currentsection, hideallsubsections]
      \end{column}
    \end{columns}
    \vfill
  \end{frame}
  \endgroup
}

% ---------- equations ----------
% Display maths is numbered, like tables and figures. The Lua filter turns
% every display equation into an `equation` environment; these two lines set
% the tag in the palette rather than in plain black, and stop beamer resetting
% the counter at each section.
% (this file is already inside \makeatletter, so @ is a letter here)
\renewcommand{\theequation}{\arabic{equation}}
\renewcommand{\tagform@}[1]{%
  \maketag@@@{\bfseries\color{dgGreen}(\ignorespaces #1\unskip\@@italiccorr)}%
}

% ---------- code ----------
% Pandoc puts a highlighted block in `Shaded` and a plain one in `verbatim`.
% Both are set small, on the same wash ground the callout boxes use, so a
% syntax example on a slide reads as another panel rather than as raw output.
% fancyvrb arrives with pandoc's highlighting, not with beamer, so a plain
% LaTeX deck that never shows code does not have \fvset at all.
\@ifundefined{fvset}{}{\fvset{fontsize=\scriptsize}}
\AtBeginDocument{%
  \@ifundefined{Shaded}{}{%
    \renewenvironment{Shaded}
      {\begin{tcolorbox}[dgbase, colback=dgWash,
          borderline west={3pt}{0pt}{dgRule},
          top=0.55em, bottom=0.55em, left=0.9em, right=0.9em]}
      {\end{tcolorbox}}%
  }%
}
\AtBeginEnvironment{verbatim}{\scriptsize}

% Pandoc ships a syntax theme in colours that are not ours. Every token is
% pulled back onto the palette: ink for code, muted for comments, blue for
% strings and keywords. Code on a slide is read, not parsed.
\newcommand{\dg@tokcolour}[2]{%
  \expandafter\ifx\csname #1\endcsname\relax\else
    \expandafter\renewcommand\csname #1\endcsname[1]{\textcolor{#2}{##1}}%
  \fi
}
\AtBeginDocument{%
  \@for\dg@tok:={DataTypeTok,DecValTok,BaseNTok,FloatTok,ConstantTok,%
    CharTok,SpecialCharTok,OtherTok,FunctionTok,VariableTok,OperatorTok,%
    BuiltInTok,ExtensionTok,AttributeTok,RegionMarkerTok,NormalTok,%
    PreprocessorTok,InformationTok,WarningTok,ErrorTok,AlertTok}\do{%
    \expandafter\dg@tokcolour\expandafter{\dg@tok}{dgInk}%
  }%
  \@for\dg@tok:={CommentTok,DocumentationTok,AnnotationTok,CommentVarTok}\do{%
    \expandafter\dg@tokcolour\expandafter{\dg@tok}{dgMuted}%
  }%
  \@for\dg@tok:={StringTok,VerbatimStringTok,SpecialStringTok,KeywordTok,%
    ControlFlowTok,ImportTok}\do{%
    \expandafter\dg@tokcolour\expandafter{\dg@tok}{dgBlue}%
  }%
}

% ---------- lists ----------
% One colour scale for both kinds of list, so depth reads the same whether
% a list is numbered or not: green, then blue, then muted. All three use the
% palette green itself -- there is no second, darker green.
% Beamer right-aligns the label in its box, so a narrow square sits further
% right than a numeral of the same level. Each marker is set in a box as wide
% as the numeral it has to line up with, flush left, so a bulleted list and a
% numbered list start on the same vertical.
% Each marker is set in a box exactly as wide as the numeral at that level,
% flush left, so the measurement maintains itself if the fonts change.
\newlength{\dg@lblw}
\newcommand{\dg@marker}[2]{%
  \settowidth{\dg@lblw}{\bfseries #1}%
  \makebox[\dg@lblw][l]{#2}%
}
\setbeamertemplate{itemize item}{%
  \dg@marker{1.}{\color{dgGreen}\rule[0.35ex]{0.5em}{0.5em}}}
\setbeamertemplate{itemize subitem}{%
  \dg@marker{a.}{\color{dgBlue}\rule[0.35ex]{0.4em}{0.4em}}}
\setbeamertemplate{itemize subsubitem}{%
  \dg@marker{i.}{\color{dgMuted}\textendash}}
\setbeamertemplate{enumerate item}{\color{dgGreen}\bfseries\insertenumlabel.}
\setbeamertemplate{enumerate subitem}{\color{dgBlue}\insertsubenumlabel.}
\setbeamertemplate{enumerate subsubitem}{\color{dgMuted}\insertsubsubenumlabel.}
\setbeamercolor{itemize/enumerate body}{fg=dgInk}
\setlength{\leftmargini}{1.4em}

% ---------- speaker notes ----------
% Hidden by default, so a ::: {.notes} block never reaches a student deck.
% The `presenter` profile includes _shared/presenter-notes.tex after this
% file and overrides it with "show notes on second screen=right".
\setbeameroption{hide notes}

% ---------- list spacing ----------
% Beamer sets list items almost on top of one another, which reads as a wall
% of text from the back of a lecture theatre. Top-level items get room to
% breathe; nested items get less than half of it, so the hierarchy stays
% visible without the block coming apart.
%
% This has to be done by redefining \@listi and friends. Setting \itemsep
% from inside the environment does not work: etoolbox's \AtBeginEnvironment
% fires before \list runs, and \list then overwrites the length with the
% value from \@listi. Tested both ways; only this one moves the type.
\def\@listi{\leftmargin\leftmargini
  \topsep 0.40em \parsep \z@ \partopsep \z@ \itemsep 0.75em}
\let\@listI\@listi
\def\@listii{\leftmargin\leftmarginii
  \labelwidth\leftmarginii \advance\labelwidth-\labelsep
  \topsep 0.25em \parsep \z@ \itemsep 0.30em}
\def\@listiii{\leftmargin\leftmarginiii
  \labelwidth\leftmarginiii \advance\labelwidth-\labelsep
  \topsep 0.20em \parsep \z@ \itemsep 0.25em}

% ---------- pandoc's \tightlist ----------
% Pandoc puts \tightlist as the first thing inside every list whose items are
% not separated by blank lines in the markdown, and defines it as
%   \setlength{\itemsep}{0pt}\setlength{\parskip}{0pt}
% That runs INSIDE the list, after \list has already read \@listi, so it
% silently overrode everything set above: measured on the Lecture 1.1 render,
% the gap between one bullet and the next was 13.5pt — exactly the line
% spacing within a single bullet, i.e. no separation at all. Only lists
% written with blank lines between items escaped it, which is why those look
% right and ordinary ones do not.
%
% Neutralising it hands the spacing back to \@listi above, so the 0.75em
% actually reaches the slide. Lists written with blank lines are unchanged:
% pandoc does not emit \tightlist for those.
%
% \def rather than \renewcommand, because this file is also usable outside
% pandoc, where \tightlist does not exist. \AtBeginDocument, because
% pandoc's own \providecommand can come after this file in the preamble.
\AtBeginDocument{\def\tightlist{}}

% ---------- blocks ----------
% Flat, with a left accent bar instead of a title bar.
\setbeamercolor{block title}{fg=dgNavy, bg=dgWash}
\setbeamercolor{block body}{fg=dgInk, bg=dgWash}
\setbeamercolor{block title alerted}{fg=white, bg=dgGreen}
\setbeamercolor{block body alerted}{fg=dgInk, bg=dgWash}
\setbeamercolor{block title example}{fg=white, bg=dgBlue}
\setbeamercolor{block body example}{fg=dgInk, bg=dgWash}
\setbeamerfont{block title}{series=\bfseries, size=\small}

% Continued frames (e.g. a long reference list) keep the plain title.
\setbeamertemplate{frametitle continuation}{}

% ---------- tables ----------
% Panelled summary/results tables, full width, booktabs rules, notes block.
%
%   \begin{dgtable}{l rrrrrr}
%     \dgpanel{7}{Panel A: County Outcome Variables}
%     \thead{Variable} & \thead{Mean} & ... \\ \midrule
%     Personal Income Growth (\%) & 2.145 & ... \\
%   \end{dgtable}
%   \dgnotes{Summary statistics for ...}
%
% Everything is black by default; use \hl{} or \hlg{} to mark the one row
% or cell that carries the argument. Rules span the full text width
% because the table is a tabular* stretched with \extracolsep.
\setlength{\aboverulesep}{0pt}\setlength{\belowrulesep}{0pt}
\setlength{\heavyrulewidth}{0.9pt}\setlength{\lightrulewidth}{0.4pt}
\arrayrulecolor{dgInk}
\renewcommand{\arraystretch}{1.18}

% Every table is \footnotesize, however it was written. \dgtable sets its own
% size, but a markdown table goes through pandoc as a longtable or a plain
% tabular and would otherwise come out at full body size -- far too big for a
% slide, since even a small econ table has more columns than body text does.
% An environment is a group, so each hook undoes itself on \end.
\AtBeginEnvironment{longtable}{\footnotesize}
\AtBeginEnvironment{tabular}{\footnotesize}
\AtBeginEnvironment{tabularx}{\footnotesize}
\AtBeginEnvironment{table}{\footnotesize}

\newcommand{\thead}[1]{\textcolor{dgInk}{\bfseries #1}}
% Colour in a table should mean something. \hl marks the row or cell that
% carries the argument; everything else stays black.
\newcommand{\hl}[1]{\textcolor{dgBlue}{\bfseries #1}}
\newcommand{\hlg}[1]{\textcolor{dgGreen}{\bfseries #1}}
% A panel heading. Emits the heading row only -- put a \midrule above it (or
% below it, if the panel opens with its own column headings).
\newcommand{\dgpanel}[2]{%
  \multicolumn{#1}{@{}l}{\textcolor{dgInk}{\bfseries #2}}\\
}
% \begin{dgtable}[\scriptsize]{l rrr} -- optional first argument sets the
% type size, for a table with more rows than \footnotesize will hold.
\newenvironment{dgtable}[2][\footnotesize]{%
  #1%
  \setlength{\tabcolsep}{0pt}%
  \begin{tabular*}{\linewidth}{@{\extracolsep{\fill}}#2@{}}
  \toprule
}{%
  \bottomrule
  \end{tabular*}%
}
% A colour chip, for the brand-reference slide.
\newcommand{\swatch}[1]{%
  \tikz[baseline=-0.35ex]{\fill[#1] (0,0) rectangle (0.75em,0.75em);}\,%
}

% Fixed-effects rows: a check rather than "Yes". Black, like the rest of the
% table -- colour is reserved for \hl and \hlg.
\newcommand{\yes}{\checkmark}
\newcommand{\no}{}

% \dgnotes lives in dublin-components.tex now, with the other callouts,
% because it is one of them: same wash panel, same left bar. It used to be
% a bold italic "Notes:" label floating in the body, which read as a
% typesetting accident next to a .defn or .takeaway box.

% ---------- captions ----------
% No "Figure 1:" label on a slide; the caption is the sentence under the
% exhibit, set muted and one size down, left aligned to the figure.
\setbeamertemplate{caption}[numbered]
\setbeamertemplate{caption label separator}{: }
\setbeamerfont{caption}{size=\footnotesize}
\setbeamerfont{caption name}{size=\footnotesize, series=\bfseries}
\setbeamercolor{caption}{fg=dgInk}
\setbeamercolor{caption name}{fg=dgInk}
% Pandoc loads the caption package for markdown tables. Left-align those
% captions at the text margin, so a markdown table's caption sits where
% \tabcaption puts one rather than centred over a narrower table.
\AtBeginDocument{%
  \@ifpackageloaded{caption}{%
    \captionsetup{justification=raggedright, singlelinecheck=false}%
  }{}%
}

% Pandoc emits \includegraphics[...,height=\textheight] for a markdown
% image, which fills the frame and pushes the caption off the slide.
% A figure environment is a group, so shrinking \textheight inside it
% bounds the image and restores itself on exit.
\AtBeginEnvironment{figure}{\setlength{\textheight}{0.50\textheight}}

% ---------- links ----------
\hypersetup{colorlinks=true, linkcolor=dgBlue, urlcolor=dgBlue,
            citecolor=dgBlue, filecolor=dgBlue}

% ---------- bibliography ----------
\setbeamertemplate{bibliography item}{}
\setbeamercolor{bibliography entry author}{fg=dgNavy}
\setbeamercolor{bibliography entry title}{fg=dgInk}
\setbeamerfont{bibliography entry title}{size=\footnotesize}
\setbeamerfont{bibliography entry author}{size=\footnotesize}

\makeatother
